%%%%%%%%%%%% STRUCTURE %%%%%%%%%%%%%%%
\documentclass[a4paper,12pt]{article}
\usepackage[T1]{fontenc}
\usepackage[utf8]{inputenc}
\usepackage[brazil]{babel}
\usepackage{lmodern}
\usepackage{setspace}
\usepackage[top=2cm, bottom=2cm, left=2cm, right=2cm]{geometry}
\usepackage{listings}
\usepackage{xcolor}
\usepackage{amsmath}
\usepackage{ragged2e}
\usepackage{float}

\lstset{
    basicstyle=\ttfamily\small,
    keywordstyle=\color{blue}\bfseries,
    stringstyle=\color{red},
    commentstyle=\color{gray}\itshape,
    numbers=left,
    numberstyle=\tiny\color{gray},
    stepnumber=1,
    numbersep=8pt,
    backgroundcolor=\color{white},
    showspaces=false,
    showstringspaces=false,
    showtabs=false,
    frame=single,
    rulecolor=\color{black!20},
    tabsize=4,
    captionpos=b,
    breaklines=true,
    breakatwhitespace=false,
}
%%%%%%%%%%%%%%%%%%%%%%%%%%%%%%%%%%%%%%

%%%%%%%%%%%%%%%% PAGES STYLE %%%%%%%%%
\usepackage{fancyhdr}
\setlength{\headheight}{15pt}
\fancypagestyle{main}{
\renewcommand{\headrulewidth}{0pt}
\fancyhead[RO]{\thepage}
\fancyfoot[CO]{}
}
%%%%%%%%%%%%%%%%%%%%%%%%%%%%%%%%%%%%%%

\usepackage{graphicx}
\usepackage{epstopdf}
\usepackage{subfig}
\usepackage{mathptmx}
\usepackage{changepage}
\usepackage{xstring}
%\usepackage[alf]{abntex2cite}

%%%%%%%%%%% PDF METADATA %%%%%%%%%%%%%
\usepackage[ pdftitle={RELATÓRIO SISTEMAS DE CONTROLE},
pdfsubject={LABORATÓRIO DE CONTROLE},
pdfkeywords={Controle,Automação,UFRN,DCA},
hidelinks]{hyperref}
%%%%%%%%%%%%%%%%%%%%%%%%%%%%%%%%%%%%%%

%%%%%%%%%%% MACROS DO TEMPLATE %%%%%%%%
% \dado{x}      -> dado experimental a preencher/conferir (vermelho no rascunho)
% \pendente{x}  -> orientação de discussão a completar (azul no rascunho,
%                  some na versão final)
% Para gerar a versão final, troque \rascunhotrue por \rascunhofalse.
\newif\ifrascunho
\rascunhotrue
\ifrascunho
  \newcommand{\dado}[1]{\textcolor{red}{\textbf{[#1]}}}
  \newcommand{\pendente}[1]{\textcolor{blue}{\textit{[PENDENTE: #1]}}}
\else
  \newcommand{\dado}[1]{#1}
  \newcommand{\pendente}[1]{}
\fi
\newcommand{\pv}{\dado{?}}

% \grafico[largura]{caminho}
% Busca automaticamente o arquivo em .pdf ou .png (mesmo se passado com ou sem .svg).
% Se a imagem ainda nao foi convertida/gerada, desenha uma caixa indicativa de pendencia.
\newcommand{\grafico}[2][0.48\textwidth]{%
  \StrSubstitute{#2}{.svg}{}[\caminholimpo]%
  \IfFileExists{\caminholimpo.pdf}{%
    \includegraphics[width=#1]{\caminholimpo.pdf}%
  }{%
    \IfFileExists{\caminholimpo.png}{%
      \includegraphics[width=#1]{\caminholimpo.png}%
    }{%
      \IfFileExists{#2}{%
        \includegraphics[width=#1]{#2}%
      }{%
        \fbox{\parbox[c][\dimexpr #1*55/100\relax][c]{\dimexpr #1-2\fboxsep-2\fboxrule\relax}{%
          \centering\footnotesize\color{red}Figura pendente\\[2pt]\texttt{\detokenize{#2}}}}%
      }%
    }%
  }%
}

% \figuraensaio{legenda}{rótulo}{saída}{sinal de controle}{erro}
% Usado para o primeiro ensaio: apresenta saída (PV e SP), sinal de controle (MV) e erro.
\newcommand{\figuraensaio}[5]{%
\begin{figure}[H]
  \centering
  \subfloat[Nível medido (PV) e referência (SP)]{\grafico[0.48\textwidth]{#3}}\hfill
  \subfloat[Sinal de controle (MV)]{\grafico[0.48\textwidth]{#4}}\\[1ex]
  \subfloat[Sinal de erro]{\grafico[0.52\textwidth]{#5}}
  \caption{#1}
  \label{#2}
\end{figure}}

% \figurasaida{legenda}{rótulo}{saída}
% Usado para os demais ensaios: apresenta somente a curva de saída (PV e SP).
\newcommand{\figurasaida}[3]{%
\begin{figure}[H]
  \centering
  \grafico[0.72\textwidth]{#3}
  \caption{#1}
  \label{#2}
\end{figure}}

% \figuracascata{legenda}{rótulo}{L2 e ref.}{L1 e L1d}{Vp}{erros}
% Figura 2x2 padrão de um ensaio com controle em cascata.
\newcommand{\figuracascata}[6]{%
\begin{figure}[H]
  \centering
  \subfloat[Nível $L_2$ e referência $r$]{\grafico{#3}}\hfill
  \subfloat[Nível $L_1$ e referência interna $L_{1d}$]{\grafico{#4}}\\
  \subfloat[Tensão na bomba $V_P$ (MV)]{\grafico{#5}}\hfill
  \subfloat[Erros $e_x = r - L_2$ e $e_i = L_{1d} - L_1$]{\grafico{#6}}
  \caption{#1}
  \label{#2}
\end{figure}}
%%%%%%%%%%%%%%%%%%%%%%%%%%%%%%%%%%%%%%

\begin{document}

\onehalfspacing

\thispagestyle{empty}

\setcounter{page}{1}

%%%%%%%%%%%% LOGOS %%%%%%%%%%%%%%%%%%%

\begin{figure}[!ht]

\centering

\subfloat{
\includegraphics[width=2.7cm]{img/UFRN.eps}
\label{UFRN Logo}
}
\hspace{11.09cm}
\subfloat{
\includegraphics[width=2.4cm]{img/DCA.eps}
\label{DCA Logo}
}

%\caption{}
\label{Logos}

\end{figure}

%%%%%%%%%%%%%%% CAPA %%%%%%%%%%%%%%%%%

\vspace{-1cm}

\begin{center}
{\bf{\normalsize UNIVERSIDADE FEDERAL DO RIO GRANDE DO NORTE\\
CENTRO DE TECNOLOGIA\\
DEPARTAMENTO DE ENGENHARIA DE COMPUTAÇÃO E AUTOMAÇÃO\\
ENGENHARIA DE COMPUTAÇÃO E ENGENHARIA MECATRÔNICA
}}

\vspace{2.8cm}

{\bf{\large RELATÓRIO DA 2ª UNIDADE\\
PROJETO DE CONTROLADORES PID\\
}}
\vspace{1.5cm}
{\large TURMA: 1\\
	GRUPO 20262G4T1}

\vspace{2.6cm}

\begin{flushright}
\begin{normalsize}
Danilo Isamu Pereira Hikague: 20230072008\\
\vspace{0.8cm}
Murilo de Lima Barros: 20250070735\\
\vspace{0.8cm}
Rafael Cabral Melo: 20210080441\\
\vspace{0.8cm}
Ramon Vinicius Ferreira de Souza: 20220017441\\
\vspace{0.8cm}
Virna Emanuelle Veloso Aguiar: 20220025354\\
\end{normalsize}
\end{flushright}

\vspace{1.5cm}

{\large Natal-RN\\
2026}

\end{center}

\newpage

%%%%%%%%%%%%%%%  CONTRA-CAPA %%%%%%%%%

\thispagestyle{empty}

\begin{center}
\begin{normalsize}
Danilo Isamu Pereira Hikague: 20230072008\\
\vspace{0.8cm}
Murilo de Lima Barros: 20250070735\\
\vspace{0.8cm}
Rafael Cabral Melo: 20210080441\\
\vspace{0.8cm}
Ramon Vinicius Ferreira de Souza: 20220017441\\
\vspace{0.8cm}
Virna Emanuelle Veloso Aguiar: 20220025354\\

\end{normalsize}
\end{center}
\vspace{3cm}

{\bf{\large {\centering PROJETO DE CONTROLADORES PID\\
}}}

\vspace{4cm}

\begin{adjustwidth}{7.5cm}{0cm}

{\normalsize

Segundo Relatório Parcial apresentado à disciplina de
Projeto de Sistemas de Controle - Laboratório, correspondente à
avaliação da 2ª unidade do semestre 2026.2 do 7º período
dos cursos: Engenharia de Computação e Engenharia Mecatrônica da
Universidade Federal do Rio Grande do Norte, sob
orientação do {\bf Prof. Fábio Meneghetti Ugulino de
Araújo.}

}

\end{adjustwidth}

\vspace{2cm}

\begin{center}

Professor:  Fábio Meneghetti Ugulino de Araújo.

\vspace{2.5cm}

{\large Natal-RN\\
2026}

\end{center}

\newpage

%%%%%%%%%%%%%%%  RESUMO %%%%%%%%%%%%%%

\thispagestyle{empty}

\begin{center}
{\large \textbf{RESUMO}}
\end{center}

\vspace{3cm}

\vspace{\baselineskip}

\begin{justify}
    \hspace{4ex}Este relatório apresenta o projeto, a implementação e a avaliação experimental de controladores P, PI, PD e PID, incluindo filtro na ação derivativa e mecanismo \textit{anti-reset windup}, aplicados ao sistema de tanques acoplados da Quanser\textsuperscript{\textregistered}. Os controladores foram implementados no ambiente MATLAB\textsuperscript{\textregistered}/SIMULINK\textsuperscript{\textregistered} com o QUARC\textsuperscript{\textregistered}, a partir de blocos básicos (ganhos, somadores, integradores e derivadores), com o sinal de controle saturado em $\pm 4$~V. Foram realizados três experimentos: o controle do nível do tanque 1, que se comporta como um sistema de primeira ordem (Experimento 2A); o controle do nível do tanque 2 na configuração 2, que se comporta como um sistema de segunda ordem, com um sistema de intertravamentos para garantir a segurança da operação (Experimento 2B); e o controle em cascata do nível do tanque 2, com uma malha interna para o tanque 1 (Experimento 2C). Para cada lei de controle, investigou-se a influência dos ganhos sobre o sobressinal, os tempos de subida e de acomodação e o erro em regime permanente. No tanque 1, os controladores PID, PID com filtro derivativo e PID com \textit{anti-reset windup} bem sintonizados levaram o nível à referência em cerca de 10~s, sem sobressinal e sem erro em regime. No tanque 2, a malha simples apresentou sobressinal e oscilações persistentes do nível do tanque 1, com acomodação de $L_2$ em, no mínimo, cerca de 55~s; o filtro derivativo eliminou o ruído da ação derivativa, e valores elevados de $K_{aw}$ esvaziaram o integrador e deixaram erro residual. O controle em cascata eliminou as oscilações e o sobressinal e reduziu o tempo de acomodação para cerca de 25~s, sendo a ação integral no controlador mestre o requisito essencial para erro nulo e o \textit{anti-reset windup} no escravo o recurso de maior impacto sobre o transitório.
\end{justify}

\vspace{\baselineskip}

\noindent\textbf{Palavras-chave:} Sistemas de controle. Tanques acoplados. Controle PID. Anti-windup. Controle em cascata.

\newpage

%%%%%%%%%%%%%%%  LISTA DE SÍMBOLOS %%%

% \thispagestyle{empty}
%
% \begin{center}
% {\large \textbf{LISTA DE SÍMBOLOS}}
% \end{center}
%
% \vspace{3cm}
%
% \begin{tabular}{ l l }
% A\hspace{1.5cm} & Matriz triangular superior com diagonal unitária.\\
% \phantom{a} & \phantom{a}\\
% D\hspace{1.5cm} & Matriz diagonal obtida a partir de $W^{T}W$\\
% \phantom{a} & \phantom{a}\\
% $\theta$\hspace{1.5cm} & Vetor de parâmetros.\\
% \phantom{a} & \phantom{a}\\
% $\Xi$\hspace{1.5cm} & Vetor de resíduos de modelagem.\\
% \phantom{a} & \phantom{a}\\
% d\hspace{1.5cm} & Tempo de retardo de um sistema ou tempo morto.\\
% \phantom{a} & \phantom{a}\\
% e(k)\hspace{1.5cm} & Resíduo (Erro de Estimação mais o Ruído).\\
% \end{tabular}
%
% \newpage

%%%% LISTA DE ABREVIATURAS E SIGLAS %%

% \thispagestyle{empty}
%
% \begin{center}
% {\large \textbf{LISTA DE ABREVIATURAS E SIGLAS}}
% \end{center}
%
% \vspace{3cm}
%
% \begin{tabular}{ l l }
% ARX\hspace{1.5cm} & Matriz triangular superior com diagonal unitária.\\
% ARMAX\hspace{1.5cm} & Matriz diagonal obtida a partir de $W^{T}W$\\
% NARX\hspace{1.5cm}&Vetor de parâmetros.\\
% NARMAX\hspace{1.5cm}&Vetor de resíduos de modelagem.\\
% MQ\hspace{1.5cm}&Tempo de retardo de um sistema ou tempo morto.\\
% \end{tabular}
%
% \newpage

%%%%%%%%% LISTA DE FIGURAS %%%%%%%%%%%

\thispagestyle{empty}

\begin{center}
\listoffigures
\end{center}

\newpage

%%%%%%%%% LISTA DE TABELAS %%%%%%%%%%%

\thispagestyle{empty}

\begin{center}
\listoftables
\end{center}

\newpage

%%%%%%%%%%%%%%% SUMÁRIO %%%%%%%%%%%%%%

\thispagestyle{empty}

\begin{center}
\tableofcontents
\end{center}

\newpage

%%%%%%%%%%%%%%% INTRODUÇÃO %%%%%%%%%%%

\thispagestyle{main}

\section{INTRODUÇÃO}

\begin{justify}
    \hspace{4ex}A introdução de um controlador em um sistema dinâmico visa modificar o seu comportamento, atuando sobre uma ou mais variáveis manipuladas (MV) para que as variáveis de processo (PV) satisfaçam especificações de desempenho tanto no regime transitório quanto no regime permanente. Dentre as estratégias de controle existentes, o controlador Proporcional-Integral-Derivativo (PID) e as suas variações (P, PI e PD) continuam sendo as mais empregadas na indústria, em virtude de sua simplicidade, robustez e da interpretação física direta de cada uma de suas ações.
\end{justify}

\begin{justify}
    \hspace{4ex}Esta segunda unidade da disciplina tem como foco o controle em malha fechada da planta dos tanques acoplados, explorando de forma progressiva as ações de controle e as modificações práticas necessárias para que o controlador opere adequadamente em um sistema real, sujeito a ruído de medição e a limites físicos do atuador.
\end{justify}

\begin{justify}
    \hspace{4ex}As atividades foram divididas em três experimentos. No Experimento 2A, controlou-se o nível do tanque 1, que se comporta como um sistema de primeira ordem, implementando-se sucessivamente as ações proporcional, integral e derivativa, o filtro na ação derivativa e o filtro \textit{anti-reset windup}. No Experimento 2B, o mesmo procedimento foi aplicado ao controle do nível do tanque 2 na configuração 2, em que o sistema passa a se comportar como um sistema de segunda ordem, e implementou-se um sistema de intertravamentos para garantir a segurança da operação. Por fim, no Experimento 2C, adotou-se a estratégia de controle em cascata, na qual uma malha externa (mestre) controla o nível do tanque 2 e uma malha interna (escrava) controla o nível do tanque 1.
\end{justify}

\begin{justify}
    \hspace{4ex}Em todos os experimentos, os controladores foram implementados no MATLAB\textsuperscript{\textregistered}/SIMULINK\textsuperscript{\textregistered} utilizando apenas blocos básicos, sem o uso de blocos PID prontos, e o sinal de controle foi saturado em $\pm 4$~V, respeitando os limites operacionais da bomba. O objetivo deste relatório é descrever a implementação realizada e analisar, com base nas curvas de resposta e nos sinais de controle obtidos, a diferença de comportamento do sistema com cada lei de controle e com diferentes valores de ganhos.
\end{justify}

\newpage

%%%%%%% FUNDAMENTAÇÃO TEÓRICA %%%%%%%

\section{FUNDAMENTAÇÃO TEÓRICA}

\subsection{Sistemas de Nível em Série}
\label{sec:ft_nivel}

\begin{justify}
    \hspace{4ex}Um tanque que descarrega livremente por um orifício na base pode ser aproximado, em torno de um ponto de operação, por um sistema de primeira ordem, $G_1(s) = K_1/(\tau_1 s + 1)$. Quando essa descarga alimenta um segundo tanque, o balanço de massa do segundo tanque, com as vazões dos orifícios dadas pela lei de Torricelli, fica como na Equação~\eqref{eq:ft_g2}, em que $A_2$ é a área da seção do segundo tanque, $a_1$ e $a_2$ são as áreas dos orifícios e $g$ é a aceleração da gravidade \cite{ogata2010}.
\end{justify}

\begin{equation}
    A_2 \frac{dL_2}{dt} = a_1\sqrt{2gL_1} - a_2\sqrt{2gL_2}
    \;\;\Rightarrow\;\;
    G_2(s) = \frac{\tilde{L}_2(s)}{\tilde{L}_1(s)} = \frac{K_2}{\tau_2 s + 1}, \quad \tau_2 = \frac{A_2}{a_2}\sqrt{\frac{2L_{20}}{g}}, \quad K_2 = \frac{L_{20}}{L_{10}}
    \label{eq:ft_g2}
\end{equation}

\begin{justify}
    \hspace{4ex}O conjunto dos dois tanques é descrito por $G_1(s)G_2(s)$, um sistema de segunda ordem com dois polos reais, mais lento que cada tanque isolado e com mais atraso de fase. Em regime, as vazões dos dois orifícios se igualam e $L_{20}/L_{10} = (a_1/a_2)^2$; com orifícios iguais, os dois tanques param no mesmo nível. Como as constantes de tempo dependem de $\sqrt{L_0}$, o comportamento muda com o ponto de operação. Esses sistemas também são do tipo 0, ou seja, não têm polo na origem, e por isso só mantêm um nível constante se receberem uma entrada de equilíbrio diferente de zero, suficiente para repor o que sai pelo orifício.
\end{justify}

\subsection{Controlador PID}

\begin{justify}
    \hspace{4ex}O controlador PID calcula o sinal de controle a partir do erro $e(t) = r(t) - y(t)$ entre a referência e a medição. Na forma paralela, cada ação tem seu próprio ganho, como mostra a Equação~\eqref{eq:ft_pid} \cite{ogata2010, astrom2006}.
\end{justify}

\begin{equation}
    u(t) = K_p\, e(t) + K_i \int_0^t e(\tau)\, d\tau + K_d \frac{de(t)}{dt}
    \label{eq:ft_pid}
\end{equation}

\begin{justify}
    \hspace{4ex}A ação proporcional responde ao erro do instante atual. Aumentar $K_p$ deixa a resposta mais rápida e diminui o erro em regime, mas não o zera: se a planta precisa de uma entrada de equilíbrio $u_0$ para se manter parada, o controlador P só a fornece se houver um erro residual $e_{ss} = u_0/K_p$. Em plantas de segunda ordem, ganhos proporcionais altos também tornam a resposta mais oscilatória.
\end{justify}

\begin{justify}
    \hspace{4ex}Esse efeito pode ser quantificado para os dois tanques em série. Com $G_1(s)G_2(s) = K/[(\tau_1 s + 1)(\tau_2 s + 1)]$, em que $K = K_1 K_2$, e um controlador P, a equação característica da malha fechada é $\tau_1\tau_2 s^2 + (\tau_1 + \tau_2)s + (1 + K K_p) = 0$. Comparando-a com a forma padrão $s^2 + 2\xi\omega_n s + \omega_n^2$, obtêm-se a frequência natural e o fator de amortecimento da Equação~\eqref{eq:xi_p}.
\end{justify}

\begin{equation}
    \omega_n = \sqrt{\frac{1 + K K_p}{\tau_1 \tau_2}}, \qquad
    \xi = \frac{\tau_1 + \tau_2}{2\sqrt{\tau_1 \tau_2\,(1 + K K_p)}}
    \label{eq:xi_p}
\end{equation}

\begin{justify}
    \hspace{4ex}Aumentar $K_p$ aumenta $\omega_n$, o que acelera a resposta, mas reduz $\xi$ na proporção de $1/\sqrt{1 + K K_p}$: a partir de $\xi < 1$, os polos tornam-se complexos conjugados e surgem sobressinal e oscilações. Ao mesmo tempo, o erro em regime para um degrau de amplitude $r$, em variáveis de desvio, vale $e_{ss} = r/(1 + K K_p)$, que diminui com $K_p$ sem nunca se anular. Há, portanto, um compromisso entre precisão em regime e amortecimento.
\end{justify}

\begin{justify}
    \hspace{4ex}A ação integral soma o erro ao longo do tempo e acrescenta um polo na origem à malha aberta. Assim, o controlador continua fornecendo a entrada de equilíbrio mesmo quando o erro chega a zero, e o erro em regime para uma referência em degrau desaparece. O preço é um atraso de fase de $90^\circ$, que aumenta o sobressinal e as oscilações, tanto mais quanto maior for $K_i$ e a ordem da planta.
\end{justify}

\begin{justify}
    \hspace{4ex}A ação derivativa reage à velocidade com que o erro varia e, com isso, antecipa a sua tendência. Ela adiciona avanço de fase e amortecimento, mas não interfere no regime permanente, já que a derivada do erro tende a zero. Em plantas de primeira ordem, que têm pouco atraso de fase, o ganho com a derivada é pequeno; em plantas de ordem maior, ele passa a ser relevante. O problema é que o ganho do derivador cresce com a frequência, de modo que o ruído de medição é amplificado e cada degrau na referência provoca um pico no sinal de controle.
\end{justify}

\subsection{Filtro na Ação Derivativa}

\begin{justify}
    \hspace{4ex}Como o derivador ideal não pode ser implementado e amplifica as altas frequências, ele costuma ser associado a um filtro passa-baixa de primeira ordem, como na Equação~\eqref{eq:ft_filtro} \cite{astrom2006}.
\end{justify}

\begin{equation}
    D(s) = \frac{K_d\, N\, s}{s + N} = \frac{K_d\, s}{s/N + 1}
    \label{eq:ft_filtro}
\end{equation}

\begin{justify}
    \hspace{4ex}Nessa forma, $N$ é a frequência de corte do filtro, em rad/s, e equivale ao inverso de sua constante de tempo, $N = 1/\tau_f$. Abaixo de $N$, o termo continua agindo como derivador; acima, seu ganho fica limitado a $K_d N$. Um $N$ grande preserva o avanço de fase, mas deixa passar mais ruído, enquanto um $N$ pequeno filtra melhor o ruído e, junto com ele, parte do amortecimento que a derivada deveria trazer. Essa é também a forma implementada no SIMULINK, em que o ramo derivativo é construído com um integrador realimentado pelo ganho $N$ (Figura~\ref{fig:2a-pidf}).
\end{justify}

\subsection{Saturação do Atuador e \textit{Windup}}

\begin{justify}
    \hspace{4ex}Todo atuador tem limites. Quando ele satura, ou quando alguma proteção bloqueia o sinal de controle, a planta deixa de receber o que o controlador calcula e a malha fica, na prática, aberta. Se o erro não muda de sinal durante esse período, o integrador continua acumulando e chega a valores bem maiores que os necessários. Quando o erro finalmente se inverte, esse excesso precisa ser descarregado antes que o sinal de controle saia da saturação, e o resultado é um sobressinal grande e uma acomodação lenta. Esse efeito é conhecido como \textit{integrator windup} \cite{astrom2006}.
\end{justify}

\begin{justify}
    \hspace{4ex}Uma forma comum de evitá-lo é o recálculo (\textit{back-calculation}), em que a diferença entre o sinal saturado $u_{sat}$ e o sinal calculado $u$ volta para a entrada do integrador multiplicada por um ganho $K_{aw}$, como na Equação~\eqref{eq:ft_aw}.
\end{justify}

\begin{equation}
    \frac{dI(t)}{dt} = K_i\, e(t) + K_{aw}\left[u_{sat}(t) - u(t)\right]
    \label{eq:ft_aw}
\end{equation}

\begin{justify}
    \hspace{4ex}Fora da saturação, esse termo é nulo e o controlador funciona como um PID comum. Durante a saturação, ele tem sinal oposto ao do erro acumulado e impede que o integrador cresça além do que o atuador consegue entregar. O valor de $K_{aw}$ define a velocidade dessa descarga: se for baixo demais, a proteção quase não atua; se for alto demais, o integrador é esvaziado antes da hora e não acumula o que o regime permanente exige.
\end{justify}

\subsection{Controle em Cascata}

\begin{justify}
    \hspace{4ex}No controle em cascata, duas malhas são aninhadas para controlar uma única variável. O controlador mestre, na malha externa, compara a variável principal com a sua referência e define a referência da malha interna, na qual o controlador escravo atua diretamente sobre o atuador a partir de uma variável intermediária \cite{seborg2011}. Com a malha interna fechada, o mestre passa a enxergar a planta da Equação~\eqref{eq:ft_cascata}, em que $C_E(s)$ é o controlador escravo, $G_1(s)$ é o processo da malha interna e $G_2(s)$ é o restante do processo.
\end{justify}

\begin{equation}
    G_{eq}(s) = T_i(s)\, G_2(s), \qquad T_i(s) = \frac{C_E(s) G_1(s)}{1 + C_E(s) G_1(s)}
    \label{eq:ft_cascata}
\end{equation}

\begin{justify}
    \hspace{4ex}Quando a malha interna é bem mais rápida que a externa, $T_i(s) \approx 1$ e o mestre enxerga praticamente só $G_2(s)$, uma planta mais simples e com menos atraso de fase. Por isso a malha interna é sintonizada primeiro, e deve ser várias vezes mais rápida que a externa. Além disso, as perturbações que entram na malha interna são corrigidas antes de chegar à variável principal. O mestre precisa de ação integral para zerar o erro em regime, enquanto no escravo a ação proporcional costuma ser suficiente, já que o integrador do mestre compensa o erro que sobra na malha interna \cite{seborg2011}.
\end{justify}

\subsection{Métricas de Desempenho}

\begin{justify}
    \hspace{4ex}As respostas ao degrau são comparadas por métricas no domínio do tempo \cite{ogata2010}. O sobressinal máximo é calculado pela Equação~\eqref{eq:ft_mp}, em que $y_{max}$ e $y_\infty$ são os valores de pico e final da saída; quando o erro em regime é nulo, $y_\infty$ coincide com a referência. O tempo de subida $t_r$ é o intervalo em que a saída vai de 10\% a 90\% da variação total, o tempo de acomodação $t_s$ é o instante a partir do qual ela não sai mais de uma faixa de $\pm 2\%$ em torno do valor final, e o erro em regime permanente é $e_{ss} = r - y_\infty$.
\end{justify}

\begin{equation}
    M_p = \frac{y_{max} - y_\infty}{y_\infty} \times 100\%
    \label{eq:ft_mp}
\end{equation}

\newpage

%%%%%%%%%% METODOLOGIA %%%%%%%%%%%%%%%

\thispagestyle{main}

\section{METODOLOGIA}

\subsection{Aparato Experimental}

\begin{justify}
    \hspace{4ex}Os experimentos foram conduzidos na bancada do sistema de tanques acoplados da Quanser\textsuperscript{\textregistered}, integrada a um microcomputador com Windows e os softwares MATLAB\textsuperscript{\textregistered}/SIMULINK\textsuperscript{\textregistered} e QUARC\textsuperscript{\textregistered}. A comunicação entre a planta e o computador foi realizada por meio da placa de aquisição de dados Q8-USB e do módulo de potência VoltPAQ-X1, que amplifica a tensão enviada pelo \textit{software} por um ganho de 3 antes de aplicá-la à bomba. Os níveis dos tanques foram obtidos a partir da tensão dos sensores de pressão, convertida em centímetros pela curva de calibração obtida na primeira unidade ($6{,}25$~cm/V).
\end{justify}

\subsection{Implementação do Controlador PID}

\begin{justify}
    \hspace{4ex}O controlador foi implementado na forma paralela da Equação~\eqref{eq:ft_pid}, utilizando apenas blocos básicos do SIMULINK\textsuperscript{\textregistered} (ganhos, somadores, integrador e derivador), sem o uso de blocos PID prontos. Ao longo do relatório, o PID com filtro na ação derivativa é abreviado como PIDF, e o PID com \textit{anti-reset windup}, como PIDAW. Nos Experimentos 2A e 2B, cada lei de controle (P, PI, PD, PID, PIDF e PIDAW) foi ensaiada com pelo menos dois conjuntos de ganhos, registrando-se o nível, o sinal de controle, o erro e a contribuição de cada ação de controle; no Experimento 2C, foi registrado um ensaio para cada combinação de controladores. Para fins de praticidade e fluidez de leitura, nem todos os dados salvos são mostrados: todas as saídas são discutidas, mas os gráficos de erro e das ações de controle são omitidos na maioria dos casos.
\end{justify}

\begin{justify}
    \hspace{4ex}A implementação foi feita de forma incremental: inicialmente apenas a ação proporcional foi utilizada e, a cada etapa, novos testes foram realizados após a inclusão da ação integral, da ação derivativa, do filtro na ação derivativa e do filtro \textit{anti-reset windup}. As ações são somadas e o resultado é limitado por um bloco de saturação em $\pm 4$~V antes de ser enviado à bomba. O derivador filtrado segue a Equação~\eqref{eq:ft_filtro}, e o \textit{anti-reset windup} foi implementado pelo método de \textit{back-calculation}, conforme a Equação~\eqref{eq:ft_aw}: a diferença entre o sinal de controle saturado ($u_{sat}$) e o não saturado ($u$) é multiplicada por $K_{aw}$ e realimentada na entrada do integrador.
\end{justify}

\begin{justify}
    \hspace{4ex}As respostas foram comparadas pelas métricas definidas na fundamentação teórica (sobressinal, tempo de subida, tempo de acomodação e erro em regime), estimadas a partir dos gráficos obtidos em cada ensaio. O sinal de controle também foi avaliado qualitativamente quanto ao tempo em saturação e à presença de ruído.
\end{justify}

\subsection{Experimento 2A: Controle do Nível do Tanque 1}

\begin{justify}
    \hspace{4ex}No Experimento 2A, o sistema foi utilizado na configuração 1, em que a bomba alimenta o tanque 1 e a variável controlada é o nível $L_1$. Nessa configuração, o processo se comporta como um sistema de primeira ordem. Como referência, utilizou-se o valor de 15~cm para este experimento e para os demais presentes neste relatório. Os ensaios deste experimento tiveram duração de 40 a 50~s, o suficiente para visualizar a convergência completa do nível à referência. O sistema projetado no SIMULINK pode ser visualizado na Figura~\ref{fig:2a-pid}.
\end{justify}

\begin{figure}[H]
    \centering
    \grafico[0.9\textwidth]{img/2a/2a-pid}
    \caption{Sistema PID implementado no SIMULINK para o experimento 2A.}
    \label{fig:2a-pid}
\end{figure}

\begin{justify}
    \hspace{4ex}Para a implementação do filtro derivativo, foi feita a adição de um integrador no ramo derivativo junto com o ganho $N$ do filtro (Equação~\eqref{eq:ft_filtro}), assim como pode ser visualizado na Figura~\ref{fig:2a-pidf}.
\end{justify}

\begin{figure}[H]
    \centering
    \grafico[0.9\textwidth]{img/2a/2a-pidf}
    \caption{Sistema PID com filtro derivativo implementado no SIMULINK para o experimento 2A.}
    \label{fig:2a-pidf}
\end{figure}

\begin{justify}
    \hspace{4ex}Para o \textit{anti-windup} o procedimento se assemelha ao realizado no filtro derivativo, mas desta vez no ganho integrativo, com a adição de uma soma e de outro ganho. O sistema implementado no SIMULINK pode ser visualizado na Figura~\ref{fig:2a-pidaw}.
\end{justify}

\begin{figure}[H]
    \centering
    \grafico[0.9\textwidth]{img/2a/2a-pidaw}
    \caption{Sistema PID com filtro \textit{anti-reset windup} no SIMULINK para o experimento 2A.}
    \label{fig:2a-pidaw}
\end{figure}

\subsection{Experimento 2B: Controle do Nível do Tanque 2}
\label{sec:metodo_2b}

\begin{justify}
    \hspace{4ex}No Experimento 2B, o sistema foi utilizado na configuração 2, em que a bomba alimenta o tanque 1, cuja vazão de saída alimenta o tanque 2. A variável controlada passa a ser o nível $L_2$, e o processo se comporta como um sistema de segunda ordem formado por dois sistemas de primeira ordem em série. O mesmo procedimento incremental do Experimento 2A foi seguido para as leis de controle P, PI, PD, PID, PID com filtro derivativo e PID com \textit{anti-reset windup}. A referência em todos os ensaios foi um degrau para $L_2 = 15$~cm. Os ensaios tiveram duração de 120~s e partiram de níveis iniciais baixos nos dois tanques.
\end{justify}

\begin{justify}
    \hspace{4ex}No mecanismo de \textit{anti-reset windup}, o sinal $u_{sat}$ da Equação~\eqref{eq:ft_aw} é o sinal obtido \textbf{após} o bloco de intertravamentos (sinal efetivamente enviado à bomba).
\end{justify}

\subsubsection{Sistema de Intertravamentos}

\begin{justify}
    \hspace{4ex}Antes dos ensaios em malha fechada, foi implementado um sistema de intertravamentos para garantir a segurança da operação, seguindo as regras definidas no roteiro:
\end{justify}

\begin{itemize}
    \item a tensão enviada à bomba é limitada à faixa de $\pm 4$~V;
    \item tensões negativas (bomba sorvendo água) são bloqueadas quando o nível do tanque 1 está abaixo de $2$~cm (margem de segurança em relação à ponta da mangueira de alimentação);
    \item quando $L_1 \geq 29$~cm e o sinal de controle é positivo, a tensão é limitada a $1{,}5$~V (valor que mantém $L_1$ próximo do limite sem ultrapassá-lo);
    \item se $L_1 \geq 30$~cm ou $L_2 \geq 30$~cm, ocorre o desligamento de segurança (\textit{shutdown}).
\end{itemize}

\subsection{Experimento 2C: Controle em Cascata}

\begin{justify}
    \hspace{4ex}No Experimento 2C, mantida a configuração 2 do Experimento 2B, adotou-se a estratégia de controle em cascata descrita no roteiro. A malha externa (mestre) recebe o erro $e_x = r - L_2$ e calcula o nível desejado para o tanque 1, $L_{1d}$. A malha interna (escrava) recebe o erro $e_i = L_{1d} - L_1$ e calcula a tensão $V_P$ aplicada à bomba. Dessa forma, o controlador mestre não atua diretamente sobre a bomba: sua saída é uma referência de nível, e o controlador escravo é o responsável por converter essa referência em tensão. As duas leis de controle são dadas pelas Equações~\eqref{eq:cascata_mestre} e~\eqref{eq:cascata_escravo}, em que $C_M(s)$ e $C_E(s)$ assumem, conforme o ensaio, as formas P, PI, PID, PID com filtro derivativo ou PID com \textit{anti-reset windup} apresentadas anteriormente.
\end{justify}

\begin{equation}
    L_{1d}(s) = C_M(s)\,\left[R(s) - L_2(s)\right]
    \label{eq:cascata_mestre}
\end{equation}

\begin{equation}
    V_P(s) = C_E(s)\,\left[L_{1d}(s) - L_1(s)\right]
    \label{eq:cascata_escravo}
\end{equation}

\begin{justify}
    \hspace{4ex}Na Equação~\eqref{eq:cascata_mestre}, $R(s)$ é a referência de nível para o tanque 2, $L_2(s)$ é o nível medido nesse tanque e o termo entre colchetes é o erro da malha externa, $e_x$. O controlador mestre $C_M(s)$ transforma esse erro em $L_{1d}(s)$, o nível que o tanque 1 deveria ter para levar $L_2$ até a referência; por isso, a saída do mestre é dada em centímetros, e não em volts. Na Equação~\eqref{eq:cascata_escravo}, o termo entre colchetes é o erro da malha interna, $e_i$, isto é, a diferença entre o nível pedido pelo mestre e o nível medido no tanque 1, e o controlador escravo $C_E(s)$ converte esse erro na tensão $V_P(s)$ enviada à bomba. Assim, as duas equações estão encadeadas: a saída da primeira é a referência da segunda, e apenas o escravo atua fisicamente sobre a planta.
\end{justify}

\begin{justify}
    \hspace{4ex}Os controladores mestre e escravo foram construídos com os mesmos blocos básicos utilizados nos Experimentos 2A e 2B (ganhos, somadores, integradores e derivadores), com o filtro derivativo e o \textit{anti-reset windup} implementados conforme as Equações~\eqref{eq:ft_filtro} e~\eqref{eq:ft_aw}. A saída do controlador escravo continua limitada a $\pm 4$~V, e a saída do controlador mestre foi limitada à faixa de 0 a 28~cm, de forma que a referência interna fique abaixo do patamar de 29~cm do intertravamento e do limite de 30~cm do tanque 1. Os intertravamentos do Experimento 2B foram mantidos.
\end{justify}

\begin{justify}
    \hspace{4ex}Foram ensaiadas 18 combinações de controladores. Na malha mestre utilizaram-se os controladores P, PI e PID solicitados no roteiro e, adicionalmente, o PID com filtro derivativo (PIDF). Na malha escrava utilizaram-se os controladores P, PI, PID e PID com filtro derivativo para todos os mestres, e o PID com \textit{anti-reset windup} (PIDAW) para os mestres PID e PIDF. Ao longo do texto, as combinações são identificadas no formato \textit{mestre}/\textit{escravo}; por exemplo, PI/P indica controlador PI na malha mestre e controlador P na malha escrava. Os ganhos finais adotados em cada combinação, ajustados empiricamente, estão listados na Tabela~\ref{tab:2c_ganhos_metodologia}.
\end{justify}

\begin{table}[H]
\centering
\caption{Ganhos utilizados em cada combinação de controladores no Experimento 2C.}
\label{tab:2c_ganhos_metodologia}
\resizebox{\textwidth}{!}{%
\begin{tabular}{|c|c|c|c|c|c|c|c|c|c|c|}
\hline
\multicolumn{2}{|c|}{\textbf{Combinação}} & \multicolumn{4}{c|}{\textbf{Mestre}} & \multicolumn{5}{c|}{\textbf{Escravo}} \\ \hline
\textbf{Mestre} & \textbf{Escravo} & $\mathbf{K_p}$ & $\mathbf{K_i}$ & $\mathbf{K_d}$ & $\mathbf{N}$ & $\mathbf{K_p}$ & $\mathbf{K_i}$ & $\mathbf{K_d}$ & $\mathbf{N}$ & $\mathbf{K_{aw}}$ \\ \hline
P    & P     & 5   & --   & --   & --    & 5    & --    & --   & --    & -- \\ \hline
P    & PI    & 4   & --   & --   & --    & 5    & 0,2   & --   & --    & -- \\ \hline
P    & PID   & 5   & --   & --   & --    & 3,8  & 0,175 & 0,1  & --    & -- \\ \hline
P    & PIDF  & 5   & --   & --   & --    & 3,5  & 0,2   & 0,01 & 0,005 & -- \\ \hline
PI   & P     & 4,5 & 0,09 & --   & --    & 2,75 & --    & --   & --    & -- \\ \hline
PI   & PI    & 4   & 0,07 & --   & --    & 3    & 0,03  & --   & --    & -- \\ \hline
PI   & PID   & 4   & 0,07 & --   & --    & 3    & 0,03  & 0,01 & --    & -- \\ \hline
PI   & PIDF  & 4   & 0,07 & --   & --    & 2,75 & 0,03  & 0,01 & 0,005 & -- \\ \hline
PID  & P     & 4,2 & 0,1  & 0,05 & --    & 2,75 & --    & --   & --    & -- \\ \hline
PID  & PI    & 4   & 0,07 & 0,05 & --    & 3    & 0,03  & --   & --    & -- \\ \hline
PID  & PID   & 4   & 0,07 & 0,05 & --    & 3    & 0,03  & 0,01 & --    & -- \\ \hline
PID  & PIDF  & 4   & 0,07 & 0,05 & --    & 2,75 & 0,03  & 0,01 & 0,005 & -- \\ \hline
PID  & PIDAW & 4   & 0,07 & 0,05 & --    & 2,75 & 0,03  & 0,01 & 0,005 & 0,02 \\ \hline
PIDF & P     & 4   & 0,1  & 0,01 & 0,005 & 2,5  & --    & --   & --    & -- \\ \hline
PIDF & PI    & 4   & 0,07 & 0,01 & 0,005 & 3    & 0,03  & --   & --    & -- \\ \hline
PIDF & PID   & 4   & 0,07 & 0,05 & 0,005 & 3    & 0,03  & 0,01 & --    & -- \\ \hline
PIDF & PIDF  & 4   & 0,07 & 0,05 & 0,005 & 2,75 & 0,03  & 0,01 & 0,005 & -- \\ \hline
PIDF & PIDAW & 4   & 0,07 & 0,05 & 0,005 & 3    & 0,05  & 0,01 & --    & 0,02 \\ \hline
\end{tabular}%
}
\end{table}

\begin{justify}
    \hspace{4ex}Em todos os ensaios, aplicou-se um degrau de referência de $r = 15$~cm para o nível do tanque 2, partindo de níveis iniciais baixos (entre 1 e 6~cm em $L_2$), e registraram-se, durante 120~s, o nível $L_1$, o nível $L_2$ e a referência $r$. Essa duração, igual à do Experimento 2B e maior que a do 2A, permite observar a acomodação completa de $L_2$, cuja dinâmica é mais lenta que a de $L_1$. Os resultados foram também comparados com os obtidos com uma única malha de controle no Experimento 2B.
\end{justify}

\begin{justify}
    \hspace{4ex}As métricas de desempenho do Experimento 2C foram estimadas por inspeção visual dos gráficos do \textit{Scope} do SIMULINK, tomando como base as linhas de grade (5~cm no eixo vertical e 20~s no eixo horizontal). Os valores apresentados são, portanto, aproximados, com resolução da ordem de $0{,}5$~cm em nível e de $5$~s em tempo. Erros em regime que não podem ser distinguidos visualmente da referência foram considerados nulos. As métricas seguem as definições da fundamentação teórica, com uma ressalva: como a faixa de 2\% em torno de 15~cm (0,3~cm) é menor que a resolução da leitura visual, o tempo de acomodação foi estimado como o instante a partir do qual $L_2$ permanece visualmente estável em torno de seu valor final. O tempo de subida foi lido como o intervalo entre 10\% e 90\% da variação de $L_2$.
\end{justify}

\newpage

%%%%%%%%%% RESULTADOS %%%%%%%%%%%%%%%

\thispagestyle{main}

\section{RESULTADOS}

\subsection{Experimento 2A: Controle do Nível do Tanque 1 (Primeira Ordem)}

\begin{justify}
    \hspace{4ex}Nesta seção são apresentados os resultados experimentais obtidos para o controle de nível do tanque 1, operando na configuração 1, que é modelado como um sistema dinâmico de primeira ordem linearizado. Para o ensaio inicial com o controlador proporcional ($K_p = 2$), reporta-se o conjunto completo de sinais de interesse: a curva de saída do processo confrontada com o sinal de referência, o sinal de controle manipulado (a tensão enviada à bomba, com saturação em $\pm 4$~V) e o sinal de erro de seguimento, como mostra a Figura~\ref{fig:2a_p_2}. Para todos os demais ensaios do experimento, apresenta-se unicamente a resposta temporal de saída da planta com a respectiva referência, seguida da análise crítica do comportamento dinâmico observado.
\end{justify}

%-------------------------------------------------------------
\subsubsection{Controlador Proporcional (P)}

\begin{justify}
    \hspace{4ex}O controlador proporcional atua com uma ação puramente amplificadora do erro instantâneo, $u(t) = K_p \, e(t)$. A manutenção do nível em uma medição constante exige uma vazão não nula de bombeamento para equilibrar a perda contínua de fluido pelo orifício inferior de saída ($F_{in} = F_{out}$). Em regime permanente, essa vazão requer uma tensão estática de equilíbrio $V_{P0} > 0$. Como o controlador P não possui termo integrador, a existência de tensão na bomba demanda obrigatoriamente um erro residual finito, $e_{ss} \approx V_{P0}/K_p$. Dessa forma, a elevação do ganho $K_p$ reduz a magnitude do erro de regime e diminui a constante de tempo de malha fechada, mas satura o atuador durante transições e amplifica o ruído de quantização do sensor de nível.
\end{justify}

\begin{justify}
    \hspace{4ex}No teste com ganho $K_p = 2$, ilustrado na Figura~\ref{fig:2a_p_2}, observa-se a resposta característica de um sistema de primeira ordem sob controle puramente proporcional: ausência de sobressinal ($M_p = 0\%$) e aproximação assintótica ao patamar de equilíbrio. Na figura observa-se um tempo de subida de cerca de 10~s e um erro em regime permanente de cerca de 1~cm; como não há oscilação, o tempo de acomodação praticamente coincide com o fim da subida. O sinal de controle antes da saturação se mantém próximo de 30 por cerca de 2~s e então desce para aproximadamente 2, onde permanece até o fim do ensaio. Como se trata de um controlador puramente proporcional, os ganhos $K_i$ e $K_d$ são nulos, e os sinais dos ramos integral e derivativo também. É possível observar ainda o sinal saturado em 4~V em contraste com o sinal antes do saturador: sem essa limitação, a bomba receberia uma tensão além de sua especificação e poderia ser danificada.

\figuraensaio{Experimento 2A -- Controlador P com $K_p = 2$.}{fig:2a_p_2}%
    {img/2a/P/P_gain_2.pdf}{img/2a/P/P_gain_2_control_signals.pdf}{img/2a/P/P_gain_2_error.pdf}
\end{justify}

%-------------------------------------------------------------
\subsubsection{Controlador Proporcional-Integral (PI)}

\begin{justify}
    \hspace{4ex}A inclusão da ação integral introduz um pólo na origem do plano $s$, em decorrência disso o integrador acumula o erro temporal e fornece a tensão de sustentação necessária à bomba mesmo quando $e(t) \to 0$, anulando teoricamente o erro em regime permanente para degraus de referência. Entretanto, a malha fechada passa a ser regida por uma dinâmica de segunda ordem, passível de oscilações e sobressinal. Além disso, a presença do limite físico de saturação da bomba ($\pm 4$~V) favorece o fenômeno de \textit{integrator windup}: durante a saturação, o integrador continua acumulando erro fictício, postergando o início do recuo da atuação e aumentando o sobressinal e o tempo de estabilização.
\end{justify}

\begin{justify}
    \hspace{4ex}Com o ganho integral $K_i = 0{,}2$ (Figura~\ref{fig:2a_pi_2}), a velocidade de correção do erro é satisfatória, mas a estabilidade relativa do sistema é inaceitável. Observa-se um sobressinal de aproximadamente 33\%, seguido de uma descida monotônica e lenta do nível, que só se acomoda perto do fim do ensaio, em cerca de 45~s.
\end{justify}

\begin{justify}
    \hspace{4ex}Para o ajuste com ganho integral menor e com o mesmo ganho proporcional, $K_p = 2$ e $K_i = 0{,}03$, apresentado na Figura~\ref{fig:2a_pi_1}, observa-se a convergência progressiva da variável de processo até o valor de referência ($e_{ss} \to 0$) sem sobressinal e praticamente sem erro de regime, com tempo de acomodação muito menor que o ensaio anterior.
\end{justify}

\begin{figure}[H]
    \centering
    \subfloat[$K_i = 0{,}2$\label{fig:2a_pi_2}]{\grafico{img/2a/PI/pi_gain_kp2_ki02.pdf}}\hfill
    \subfloat[$K_i = 0{,}03$\label{fig:2a_pi_1}]{\grafico{img/2a/PI/pi_gain_kp2_ki003.pdf}}
    \caption{Experimento 2A -- Controlador PI com $K_p = 2$.}
    \label{fig:2a_pi_2_1_grp}
\end{figure}

%-------------------------------------------------------------
\subsubsection{Controlador Proporcional-Derivativo (PD)}

\begin{justify}
    \hspace{4ex}A ação derivativa introduz um efeito de antecipação à dinâmica do erro, acrescentando amortecimento à malha. Contudo, em um sistema de primeira ordem, o ganho derivativo puro não altera a ordem do sistema linear e não possui capacidade de eliminar o erro de regime em degrau, pois $\dot{e}(t) \to 0$ quando as variáveis estabilizam. Adicionalmente, o derivador puro possui ganho infinito para frequências muito altas, atuando como amplificador de ruídos provenientes do sensor e gerando impulsos bruscos de controle a cada alteração em degrau da referência.
\end{justify}

\begin{justify}
    \hspace{4ex}Na Figura~\ref{fig:2a_pd_1}, correspondente ao ensaio com $K_d = 1$, constata-se a persistência do erro em regime permanente. Observa-se um erro de regime $e_{ss} \approx 2$~cm e tempo de subida $t_r \approx 7$~s; nota-se também um ruído elevado no regime permanente, causado pela ação derivativa, que persiste até o fim do ensaio, em 50~s.
\end{justify}

\begin{justify}
    \hspace{4ex}Ao diminuir substancialmente o ganho derivativo para $K_d = 0{,}1$ (Figura~\ref{fig:2a_pd_0.1}), mantendo o mesmo ganho proporcional do teste anterior, nota-se um aumento no tempo de subida para cerca de 8~s, uma diminuição do erro de regime para aproximadamente 1~cm e uma redução significativa do ruído.
\end{justify}

\begin{figure}[H]
    \centering
    \subfloat[$K_d = 1$\label{fig:2a_pd_1}]{\grafico{img/2a/PD/pd-1-saida.pdf}}\hfill
    \subfloat[$K_d = 0{,}1$\label{fig:2a_pd_0.1}]{\grafico{img/2a/PD/pd-0.1-saida.pdf}}
    \caption{Experimento 2A -- Controlador PD com $K_p = 2$.}
    \label{fig:2a_pd_1_0.1_grp}
\end{figure}

%-------------------------------------------------------------
\subsubsection{Controlador Proporcional-Integral-Derivativo (PID)}

\begin{justify}
    \hspace{4ex}Nesta implementação ideal preliminar do PID sem filtro na derivada e sem tratamento de saturação, o derivador puro continua suscetível a ruídos do sensor de nível e a picos instantâneos a cada variação da referência.
\end{justify}

\begin{justify}
    \hspace{4ex}A resposta temporal com ganhos ($K_p = 1{,}8$, $K_i = 0{,}03$ e $K_d = 0{,}1$) é ilustrada na Figura~\ref{fig:2a_pid_1}. Observa-se um sobressinal praticamente nulo, tempo de acomodação $t_s \approx 10$~s e tempo de subida também $t_r \approx 10$~s, pois praticamente não há oscilação após a subida, muito semelhante aos resultados do PI.
\end{justify}

\begin{justify}
    \hspace{4ex}Elevando-se simultaneamente as ações integral e derivativa ($K_i = 0{,}1$ e $K_d = 0{,}5$), conforme a Figura~\ref{fig:2a_pid_2}, busca-se acelerar a resposta transitória, mas o que se obtém é um sobressinal elevado e um tempo de acomodação maior que a duração do ensaio: em 50~s, o nível do tanque 1 ainda está acima de 15~cm, embora com tendência a convergir para a referência.
\end{justify}

\begin{figure}[H]
    \centering
    \subfloat[$K_i = 0{,}03$ e $K_d = 0{,}1$\label{fig:2a_pid_1}]{\grafico{img/2a/PID/p1.8-i0.03-d0.1-saida.pdf}}\hfill
    \subfloat[$K_i = 0{,}1$ e $K_d = 0{,}5$\label{fig:2a_pid_2}]{\grafico{img/2a/PID/p1.8-i0.1-d0.5-saida.pdf}}
    \caption{Experimento 2A -- Controlador PID com $K_p = 1{,}8$.}
    \label{fig:2a_pid_1_2_grp}
\end{figure}

%-------------------------------------------------------------
\subsubsection{Controlador PID com Filtro na Ação Derivativa}

\begin{justify}
    \hspace{4ex}Com a adição do filtro contra ruído na ação derivativa, no primeiro ensaio (Figura~\ref{fig:2a_pidf_1}), com $K_p = 2$, $K_i = 0{,}08$, $K_d = 0{,}05$ e $N = 0{,}1$, verifica-se a regularização do perfil de subida e um ruído menor no regime permanente; porém, o ganho integral mais alto produz sobressinal e tempo de acomodação elevados, com o nível máximo chegando a aproximadamente 17~cm e $t_s \approx 45$~s.
\end{justify}

\begin{justify}
    \hspace{4ex}Mesmo com a tendência de convergência à referência, esta configuração se mostra pouco eficiente para o controle do nível do tanque.
\end{justify}

\begin{justify}
    \hspace{4ex}De forma semelhante ao resultado da Figura~\ref{fig:2a_pid_1}, o segundo ensaio (Figura~\ref{fig:2a_pidf_2}), com $K_p = 2$, $K_i = 0{,}03$, $K_d = 0{,}13$ e $N = 0{,}05$, apresenta sobressinal nulo, tempo de acomodação baixo e tempo de subida aceitável. A diferença de ruído em relação ao PID sem filtro é pequena, pois o sinal de controle do PID sem filtro já não apresentava ruído de regime elevado; por isso, os sinais de controle com e sem filtro são bastante semelhantes.
\end{justify}
  
\begin{figure}[H]
    \centering
    \subfloat[$K_i = 0{,}08$, $K_d = 0{,}05$, $N = 0{,}1$\label{fig:2a_pidf_1}]{\grafico{img/2a/PID-F/p2-i008-d005-f01-saida.pdf}}\hfill
    \subfloat[$K_i = 0{,}03$, $K_d = 0{,}13$, $N = 0{,}05$\label{fig:2a_pidf_2}]{\grafico{img/2a/PID-F/p2-i003-d013-f005-saida.pdf}}
    \caption{Experimento 2A -- PIDF com $K_p = 2$.}
    \label{fig:2a_pidf_1_2_grp}
\end{figure}

%-------------------------------------------------------------
\subsubsection{Controlador PID com Filtro \textit{Anti-Reset Windup}}

\begin{justify}
    \hspace{4ex}A implementação da malha de \textit{anti-reset windup} (via recálculo por realimentação da diferença entre o sinal de controle saturado e o não saturado, ponderada pelo ganho $K_{aw}$) tem como finalidade impedir o acúmulo excessivo de erro no integrador durante os períodos em que o atuador opera nos limites físicos de saturação ($\pm 4$~V). Sem essa proteção, o termo integral continua crescendo enquanto a bomba já está em sua capacidade máxima, provocando atraso na saída da região de saturação e elevado sobressinal.
\end{justify}

\begin{justify}
    \hspace{4ex}No primeiro ensaio desta estrutura, ilustrado na Figura~\ref{fig:2a_pidaw_1}, utilizou-se a configuração $K_p = 2$, $K_i = 0{,}035$, $K_d = 0{,}08$ e ganho de \textit{anti-windup} $K_{aw} = 0{,}001$. Observa-se uma resposta transitória limpa e bem amortecida, atingindo a referência de 15~cm sem a ocorrência de sobressinal ($M_p = 0\%$), com tempo de subida $t_r \approx 10$~s, tempo de acomodação $t_s \approx 10$~s e erro em regime permanente nulo. Adicionalmente, a combinação moderada do ganho derivativo com a descarga controlada da saturação resultou em um regime permanente praticamente isento de ruído.
\end{justify}

\begin{justify}
    \hspace{4ex}Por outro lado, no segundo ensaio (Figura~\ref{fig:2a_pidaw_2}), com $K_p = 2$, $K_i = 0{,}05$, $K_d = 0{,}1$ e $K_{aw} = 0{,}1$, observa-se o efeito adverso de uma ação de descarregamento excessivamente restritiva sobre o ramo integrador. Durante todo o ensaio de 50~s, a curva de nível permanece abaixo da referência de 15~cm, apresentando uma subida excessivamente lenta e incapacidade prática de eliminar o erro estático dentro do horizonte de teste. Isso ocorre porque a atuação severa da malha de \textit{anti-windup} durante a saturação inicial drena prematuramente a energia acumulada pelo integrador, impedindo que o controlador estabeleça rapidamente a tensão de sustentação necessária para vencer a vazão de saída do orifício.
\end{justify}

\begin{figure}[H]
    \centering
    \subfloat[$K_i = 0{,}035$, $K_d = 0{,}08$, $K_{aw} = 0{,}001$\label{fig:2a_pidaw_1}]{\grafico{img/2a/PID-AW/p2-i0035-d008-aw0001-saida.pdf}}\hfill
    \subfloat[$K_i = 0{,}05$, $K_d = 0{,}1$, $K_{aw} = 0{,}1$\label{fig:2a_pidaw_2}]{\grafico{img/2a/PID-AW/p2-i005-d01-aw01-saida.pdf}}
    \caption{Experimento 2A -- PIDAW com $K_p = 2$.}
    \label{fig:2a_pidaw_1_2_grp}
\end{figure}

%-------------------------------------------------------------
\subsubsection{Discussão Comparativa do Experimento 2A}

\begin{justify}
    \hspace{4ex}A análise integrada das seis estratégias de controle aplicadas ao tanque 1 evidencia as particularidades de cada ação frente a uma planta predominantemente de primeira ordem sujeita à saturação do atuador:
\end{justify}

\begin{itemize}
    \item \textbf{Ação Proporcional (P e PD):} Mostrou-se incapaz de zerar o erro em regime permanente em todos os testes ($e_{ss} \neq 0$), visto que o sistema exige um erro residual constante para gerar a tensão de equilíbrio da bomba. A adição pura do termo derivativo (PD) não solucionou o erro estático e amplificou os ruídos de alta frequência oriundos do sensor quando utilizado $K_d$ elevado.
    \item \textbf{Ação Integral (PI):} Mostrou-se indispensável para garantir erro estático nulo ($e_{ss} \to 0$). Enquanto ganhos integrais elevados ($K_i = 0{,}2$) provocaram forte sobressinal ($33\%$) e longo tempo de acomodação devido ao efeito de \textit{windup}, uma sintonia conservadora ($K_i = 0{,}03$) permitiu atingir a referência de forma suave.
    \item \textbf{Ação Derivativa e Filtragem (PID e PIDF):} Quando sintonizados com ganho derivativo baixo ($K_d = 0{,}1$ no PID e $K_d = 0{,}13$ com $N = 0{,}05$ no PIDF), ambos apresentaram respostas rápidas, sem sobressinal e com ruído desprezível em regime permanente, demonstrando que para esta faixa de ganhos o filtro derivativo adicional altera muito pouco a dinâmica global.
    \item \textbf{Compensação de Saturação (PIDAW):} Com sintonia adequada ($K_{aw} = 0{,}001$), proporcionou desempenho equivalente às melhores curvas do PID e do PIDF. Contudo, evidenciou-se que o ajuste inadequado do parâmetro de \textit{anti-windup} pode inibir excessivamente a ação integral, mantendo a resposta permanentemente abaixo da referência durante todo o ensaio.
\end{itemize}

\begin{justify}
    \hspace{4ex}Em síntese, os melhores resultados experimentais foram alcançados pelas estruturas PID ($K_p = 1{,}8$, $K_i = 0{,}03$ e $K_d = 0{,}1$), PIDF ($K_p = 2$, $K_i = 0{,}03$, $K_d = 0{,}13$ e $N = 0{,}05$) e PIDAW ($K_p = 2$, $K_i = 0{,}035$, $K_d = 0{,}08$ e $K_{aw} = 0{,}001$), que apresentaram comportamentos dinâmicos praticamente idênticos, com tempo de acomodação $t_s \approx 10$~s, sobressinal nulo ($M_p = 0\%$) e erro em regime permanente igual a zero. Pelo critério de simplicidade estrutural, menor custo computacional e menor número de parâmetros de sintonia, elege-se o controlador PID clássico com $K_p = 1{,}8$, $K_i = 0{,}03$ e $K_d = 0{,}1$ como a melhor solução prática para o controle de nível do tanque 1.
\end{justify}

\newpage

%=============================================================
\subsection{Experimento 2B: Controle do Nível do Tanque 2 (Segunda Ordem)}

\begin{justify}
    \hspace{4ex}Nesta seção são apresentados a validação do sistema de intertravamentos e os resultados do controle do nível do tanque 2 na configuração 2. Diferentemente do Experimento 2A, a planta vista pelo controlador é formada por dois tanques em série, o que adiciona um polo à malha e torna o sistema mais propenso a sobressinal e oscilações.
\end{justify}

\subsubsection{Validação do Sistema de Intertravamentos}

\begin{justify}
    \hspace{4ex}A validação confirmou a atuação dos intertravamentos nas condições testadas (Tabela~\ref{tab:2b_intertravamento}). Quando submetido a degraus elevados de referência, o nível $L_1$ atingiu $30$~cm e o intertravamento cortou a tensão aplicada à bomba, impedindo o transbordamento. Enquanto o intertravamento atua, a malha de controle fica efetivamente aberta e o termo integral tende a acumular erro, reforçando a necessidade do mecanismo \textit{anti-reset windup}.
\end{justify}

\begin{table}[H]
\centering
\caption{Validação experimental do sistema de intertravamentos de segurança.}
\label{tab:2b_intertravamento}
\resizebox{\textwidth}{!}{%
\begin{tabular}{|l|l|l|c|}
\hline
\textbf{Condição testada} & \textbf{Ação esperada} & \textbf{Ação observada} & \textbf{OK?} \\ \hline
Sinal fora da faixa $\pm 4$~V & Saturação em $\pm 4$~V & Tensão limitada em $+4{,}0$~V e $-4{,}0$~V & SIM \\ \hline
$u < 0$ com $L_1 < 2$~cm & Forçar tensão a $0$~V & Tensão mantida em $0{,}0$~V & SIM \\ \hline
$L_1 \geq 29$~cm e $u > 0$ & Manter $1{,}5$~V & Tensão limitada corretamente em $1{,}5$~V & SIM \\ \hline
$L_1 \geq 30$~cm & Desligamento de segurança & Tensão cortada para $0{,}0$~V imediatamente & SIM \\ \hline
$L_2 \geq 30$~cm & Desligamento de segurança & Tensão cortada para $0{,}0$~V imediatamente & SIM \\ \hline
\end{tabular}%
}
\end{table}

\begin{justify}
    \hspace{4ex}Nos gráficos a seguir, a curva verde corresponde ao nível do tanque 1 ($L_1$), a curva vermelha ao nível do tanque 2 ($L_2$), que é a variável controlada, e a linha horizontal em 15~cm à referência. Os valores citados foram estimados visualmente a partir dos gráficos e estão resumidos na Tabela~\ref{tab:2b_comparativo}. Em todos os ensaios, a bomba permanece saturada no início e $L_1$ sobe segundo praticamente a mesma rampa, de modo que o tempo de subida de $L_2$ ficou entre 15 e 20~s em todos os casos, limitado pelo atuador e não pela sintonia.
\end{justify}

\subsubsection{Controlador Proporcional (P)}

\begin{justify}
    \hspace{4ex}Assim como no tanque 1, a planta da configuração 2 é do tipo 0, e o controlador P mantém um erro em regime que diminui com o aumento de $K_p$. A diferença é que, por se tratar de um sistema de segunda ordem, o aumento do ganho reduz o fator de amortecimento $\xi$ (Equação~\eqref{eq:xi_p}) e, a partir de certo valor, torna os polos de malha fechada complexos conjugados: a resposta passa de sobreamortecida a subamortecida, com sobressinal e oscilações crescentes.
\end{justify}

\begin{justify}
    \hspace{4ex}No ensaio com $K_p = 1$ (Figura~\ref{fig:2b_p}), $L_2$ sobe até cerca de 15~cm aos 30~s e depois se acomoda em torno de 13~cm, deixando um erro em regime de aproximadamente 2~cm. Mesmo com um ganho baixo, a resposta já é subamortecida: $L_2$ ultrapassa seu valor final em cerca de 15\%, e $L_1$ chega a cerca de 28~cm, desce para cerca de 9~cm e oscila até se estabilizar em torno de 15~cm, após aproximadamente 70~s. Esse comportamento contrasta com o do tanque 1 no Experimento 2A, em que o controlador P produziu respostas sem sobressinal, e confirma a menor margem de amortecimento da planta de segunda ordem.
\end{justify}

\begin{figure}[H]
    \centering
    \grafico[0.6\textwidth]{img/2b/P/P_gain_1.pdf}
    \caption{Experimento 2B -- Controlador P com $K_p = 1$. Neste gráfico, a curva tracejada, identificada na legenda como referência, corresponde ao nível $L_1$, e a curva contínua, ao nível $L_2$.}
    \label{fig:2b_p}
\end{figure}

\subsubsection{Controlador Proporcional-Integral (PI)}

\begin{justify}
    \hspace{4ex}A ação integral elimina o erro em regime, mas eleva a ordem da malha fechada para três. Como o atraso de fase introduzido pelo integrador se soma ao dos dois tanques, o sistema torna-se mais oscilatório que no Experimento 2A, e o efeito de \textit{windup} é agravado pelos longos períodos de saturação da bomba.
\end{justify}

\begin{justify}
    \hspace{4ex}Com $K_p = 3{,}5$ e $K_i = 0{,}01$ (Figura~\ref{fig:2b_pi_1}), o erro em regime foi eliminado: $L_2$ chega a cerca de 16~cm (sobressinal de aproximadamente 7\%) e se acomoda em 15~cm por volta de 70~s. Contudo, $L_1$ oscila de forma pouco amortecida durante todo o ensaio, entre cerca de 13 e 22~cm no início e ainda com amplitude de cerca de 1~cm ao final, o que mostra que a malha simples tem dificuldade em amortecer a dinâmica do tanque intermediário.
\end{justify}

\begin{justify}
    \hspace{4ex}Com $K_i = 1$ (Figura~\ref{fig:2b_pi_2}), o efeito de \textit{windup} torna o sistema praticamente instável. O integrador acumula erro rapidamente e mantém a bomba saturada; $L_1$ fica preso no patamar do intertravamento, em torno de 28~cm, por cerca de 40~s, enquanto $L_2$ sobe até cerca de 23~cm, um sobressinal de aproximadamente 50\%. Quando o erro finalmente se inverte, o integrador precisa ser descarregado, a bomba é desligada e $L_1$ cai até cerca de 1~cm, fazendo $L_2$ descer até cerca de 7~cm; em seguida, o ciclo recomeça. O resultado é uma oscilação de grande amplitude sem acomodação, que ilustra a importância de limitar o ganho integral e de tratar a saturação.
\end{justify}

\begin{figure}[H]
    \centering
    \subfloat[$K_i = 0{,}01$\label{fig:2b_pi_1}]{\grafico{img/2b/PI/saida_p3.5_ki0.01.pdf}}\hfill
    \subfloat[$K_i = 1$\label{fig:2b_pi_2}]{\grafico{img/2b/PI/saida_p3.5_ki1.pdf}}
    \caption{Experimento 2B -- Controlador PI com $K_p = 3{,}5$. No gráfico (b), a curva tracejada, identificada na legenda como referência, corresponde ao nível $L_1$, e a curva contínua, ao nível $L_2$.}
    \label{fig:2b_pi}
\end{figure}

\subsubsection{Controlador Proporcional-Derivativo (PD)}

\begin{justify}
    \hspace{4ex}Foram realizados ensaios com $K_p = 2$ e $K_d = 0{,}3$ e com $K_p = 1{,}7$ e $K_d = 0{,}8$ (Figura~\ref{fig:2b_pd}). Por não possuir ação integral, o PD manteve erro em regime nos dois casos: $L_2$ se acomodou em cerca de 14~cm no primeiro ensaio e em cerca de 13,7~cm no segundo, que tem $K_p$ menor. Nos dois ensaios, $L_2$ chegou a cerca de 15~cm antes de se acomodar, um sobressinal da ordem de 10\% em relação ao valor final. A principal diferença está em $L_1$: com $K_d = 0{,}8$, as oscilações do tanque 1 são amortecidas mais rapidamente, o que evidencia o avanço de fase introduzido pela ação derivativa.
\end{justify}

\begin{figure}[H]
    \centering
    \subfloat[$K_p = 2$ e $K_d = 0{,}3$\label{fig:2b_pd_1}]{\grafico{img/2b/PD/p-2-d-0.3-saida.pdf}}\hfill
    \subfloat[$K_p = 1{,}7$ e $K_d = 0{,}8$\label{fig:2b_pd_2}]{\grafico{img/2b/PD/p-1.7-d-0.8-saida.pdf}}
    \caption{Experimento 2B -- Controlador PD.}
    \label{fig:2b_pd}
\end{figure}

\subsubsection{Controlador PID}

\begin{justify}
    \hspace{4ex}Ao contrário do sistema de primeira ordem, no sistema de segunda ordem a ação derivativa tem um papel relevante: o avanço de fase que ela introduz aumenta o amortecimento da malha fechada, permitindo reduzir o sobressinal e as oscilações provocadas pelas ações proporcional e integral.
\end{justify}

\begin{justify}
    \hspace{4ex}Com $K_p = 3$, $K_i = 0{,}05$ e $K_d = 0{,}2$ (Figura~\ref{fig:2b_pid_1}), o ganho integral relativamente alto prevaleceu sobre o amortecimento da derivada: $L_1$ permaneceu no patamar do intertravamento entre cerca de 20 e 32~s, $L_2$ chegou a cerca de 18~cm (sobressinal de aproximadamente 20\%) e os dois níveis continuaram oscilando até o fim do ensaio, sem acomodação em 120~s. Reduzindo $K_i$ para $0{,}01$ e aumentando $K_d$ para $0{,}8$, com $K_p = 3{,}5$ (Figura~\ref{fig:2b_pid_2}), o sobressinal de $L_2$ caiu para cerca de 7\% e o nível se acomodou em aproximadamente 55~s, o melhor tempo entre os ensaios de malha simples. As oscilações de $L_1$, porém, persistiram com amplitude decrescente até o fim do ensaio. Nesse ensaio, a parcela derivativa do sinal de controle apresentou ruído intenso, com picos da ordem de $\pm 5$, consequência da derivação do ruído de medição.
\end{justify}

\begin{figure}[H]
    \centering
    \subfloat[$K_p = 3$, $K_i = 0{,}05$ e $K_d = 0{,}2$\label{fig:2b_pid_1}]{\grafico{img/2b/PID/saida_p3_ki0.05_kd_0.2.pdf}}\hfill
    \subfloat[$K_p = 3{,}5$, $K_i = 0{,}01$ e $K_d = 0{,}8$\label{fig:2b_pid_2}]{\grafico{img/2b/PID/saida_p3.5_ki0.01_kd_0.8.pdf}}
    \caption{Experimento 2B -- Controlador PID.}
    \label{fig:2b_pid}
\end{figure}

\subsubsection{Controlador PID com Filtro na Ação Derivativa}

\begin{justify}
    \hspace{4ex}O filtro foi aplicado ao PID com $K_p = 3{,}5$ e $K_i = 0{,}01$, com $K_d = 0{,}8$ e $N = 0{,}1$ e com $K_d = 1{,}5$ e $N = 0{,}05$ (Figura~\ref{fig:2b_pidf}). Nos dois casos, a parcela derivativa do sinal de controle deixou de apresentar o ruído observado no PID sem filtro e passou a variar de forma suave (Figura~\ref{fig:2b_acoes}). Em $L_2$, o desempenho foi semelhante ao do melhor PID: sobressinal de cerca de 7\% e acomodação em aproximadamente 70~s com $K_d = 0{,}8$ e 60~s com $K_d = 1{,}5$.
\end{justify}

\begin{justify}
    \hspace{4ex}Por outro lado, com $K_d = 0{,}8$ as oscilações de $L_1$ ficaram maiores que no PID sem filtro (entre cerca de 9 e 22~cm), pois o filtro também atenua parte do avanço de fase que a derivada deveria fornecer. Com $K_d = 1{,}5$ e $N = 0{,}05$, as oscilações de $L_1$ ficaram restritas a cerca de 12 a 18~cm, sem o ruído do derivador puro; como os dois parâmetros foram alterados ao mesmo tempo, a melhora resulta da combinação de um ganho derivativo maior com um filtro de frequência de corte mais baixa. Por isso, esse ensaio representa o melhor compromisso entre velocidade, sobressinal e ruído entre os ensaios de malha simples.
\end{justify}

\begin{figure}[H]
    \centering
    \subfloat[$K_d = 0{,}8$ e $N = 0{,}1$\label{fig:2b_pidf_1}]{\grafico{img/2b/PID-F/saida-p3.5-i0.01-d0.8-fn0.1.pdf}}\hfill
    \subfloat[$K_d = 1{,}5$ e $N = 0{,}05$\label{fig:2b_pidf_2}]{\grafico{img/2b/PID-F/saida-p3.5-i0.01-d1.5-fn0.05.pdf}}
    \caption{Experimento 2B -- PIDF com $K_p = 3{,}5$ e $K_i = 0{,}01$.}
    \label{fig:2b_pidf}
\end{figure}

\begin{figure}[H]
    \centering
    \subfloat[PID sem filtro ($K_d = 0{,}8$)]{\grafico{img/2b/PID/actions_p3.5_ki0.01_kd_0.8.pdf}}\hfill
    \subfloat[PIDF ($K_d = 0{,}8$ e $N = 0{,}1$)]{\grafico{img/2b/PID-F/actions-p3.5-i0.01-d0.8-fn0.1.pdf}}
    \caption{Experimento 2B -- Parcelas proporcional, integral e derivativa do sinal de controle, com $K_p = 3{,}5$ e $K_i = 0{,}01$.}
    \label{fig:2b_acoes}
\end{figure}

\subsubsection{Controlador PID com Filtro \textit{Anti-Reset Windup}}

\begin{justify}
    \hspace{4ex}Devido à dinâmica mais lenta do tanque 2, os períodos de saturação do sinal de controle são mais longos que no Experimento 2A e, consequentemente, o acúmulo indevido de erro no integrador é maior. O \textit{anti-reset windup} foi aplicado ao PID com $K_p = 3{,}5$, $K_i = 0{,}01$ e $K_d = 0{,}8$, sem filtro derivativo, com três valores de $K_{aw}$ (Figura~\ref{fig:2b_pidaw}).
\end{justify}

\begin{justify}
    \hspace{4ex}Com $K_{aw} = 0{,}001$, a resposta de $L_2$ foi semelhante à do PID sem \textit{anti-windup}, com sobressinal de cerca de 7\% e acomodação em aproximadamente 70~s, mas $L_1$ deixou de oscilar contra o patamar do intertravamento no início do ensaio, subindo e descendo de forma suave. Com $K_{aw} = 0{,}005$ e $K_{aw} = 0{,}05$, o sobressinal de $L_2$ caiu para cerca de 3\%, porém o integrador foi descarregado em excesso durante a saturação: após o pico, $L_2$ desceu até cerca de 13~cm e voltou a subir muito lentamente, terminando o ensaio ainda abaixo da referência, com erro de cerca de 0,5~cm e 0,7~cm, respectivamente. Esse efeito, também observado no Experimento 2A, mostra que um $K_{aw}$ elevado impede o integrador de acumular o que o regime permanente exige. Como o filtro derivativo não foi utilizado nesses ensaios, a parcela derivativa manteve o mesmo ruído do PID sem filtro.
\end{justify}

\begin{figure}[H]
    \centering
    \subfloat[$K_{aw} = 0{,}001$\label{fig:2b_pidaw_1}]{\grafico{img/2b/PID-AW/saida-p3.5-i-0.01-d0.8-aw0.001.png}}\hfill
    \subfloat[$K_{aw} = 0{,}005$\label{fig:2b_pidaw_2}]{\grafico{img/2b/PID-AW/saida-p3.5-i-0.01-d0.8-aw0.005.pdf}}\\
    \subfloat[$K_{aw} = 0{,}05$\label{fig:2b_pidaw_3}]{\grafico{img/2b/PID-AW/saida-p3.5-i-0.01-d0.8-aw0.05.pdf}}
    \caption{Experimento 2B -- PIDAW com $K_p = 3{,}5$, $K_i = 0{,}01$ e $K_d = 0{,}8$.}
    \label{fig:2b_pidaw}
\end{figure}

\subsubsection{Comparativo entre as Leis de Controle}
\label{sec:2b_comparativo}

\begin{justify}
    \hspace{4ex}A Tabela~\ref{tab:2b_comparativo} reúne as métricas aproximadas de todos os ensaios do Experimento 2B.
\end{justify}

\begin{table}[H]
\centering
\caption{Métricas aproximadas dos ensaios do Experimento 2B, estimadas a partir dos gráficos.}
\label{tab:2b_comparativo}
\resizebox{\textwidth}{!}{%
\begin{tabular}{|c|c|c|c|c|}
\hline
\textbf{Controlador} & \textbf{Ganhos} & $\mathbf{M_p}$ \textbf{(\%)} & $\mathbf{t_s}$ \textbf{(s)} & $\mathbf{e_{ss}}$ \textbf{(cm)} \\ \hline
P     & $K_p = 1$ & $\approx 15$ & $\approx 70$ & $\approx 2$ \\ \hline
PI    & $K_p = 3{,}5$, $K_i = 0{,}01$ & $\approx 7$ & $\approx 70$ & 0 \\ \hline
PI    & $K_p = 3{,}5$, $K_i = 1$ & $\approx 50$ & não acomoda & -- \\ \hline
PD    & $K_p = 2$, $K_d = 0{,}3$ & $\approx 10$ & $\approx 70$ & $\approx 1$ \\ \hline
PD    & $K_p = 1{,}7$, $K_d = 0{,}8$ & $\approx 10$ & $\approx 60$ & $\approx 1{,}3$ \\ \hline
PID   & $K_p = 3$, $K_i = 0{,}05$, $K_d = 0{,}2$ & $\approx 20$ & $>120$ & oscila \\ \hline
PID   & $K_p = 3{,}5$, $K_i = 0{,}01$, $K_d = 0{,}8$ & $\approx 7$ & $\approx 55$ & 0 \\ \hline
PIDF  & $K_p = 3{,}5$, $K_i = 0{,}01$, $K_d = 0{,}8$, $N = 0{,}1$ & $\approx 7$ & $\approx 70$ & 0 \\ \hline
PIDF  & $K_p = 3{,}5$, $K_i = 0{,}01$, $K_d = 1{,}5$, $N = 0{,}05$ & $\approx 7$ & $\approx 60$ & 0 \\ \hline
PIDAW & $K_p = 3{,}5$, $K_i = 0{,}01$, $K_d = 0{,}8$, $K_{aw} = 0{,}001$ & $\approx 7$ & $\approx 70$ & 0 \\ \hline
PIDAW & $K_p = 3{,}5$, $K_i = 0{,}01$, $K_d = 0{,}8$, $K_{aw} = 0{,}005$ & $\approx 3$ & $>120$ & $\approx 0{,}5$ \\ \hline
PIDAW & $K_p = 3{,}5$, $K_i = 0{,}01$, $K_d = 0{,}8$, $K_{aw} = 0{,}05$ & $\approx 3$ & $>120$ & $\approx 0{,}7$ \\ \hline
\end{tabular}%
}
\end{table}

\begin{justify}
    \hspace{4ex}A análise comparativa mostra que: (i) o controle P não anula o erro em regime e, mesmo com ganho baixo, já produz resposta subamortecida; (ii) a ação integral anula o erro, mas ganhos integrais altos levam a \textit{windup} severo e a oscilações de grande amplitude; (iii) a ação derivativa reduz o sobressinal e o tempo de acomodação de $L_2$, desde que acompanhada de filtro para eliminar o ruído no sinal de controle; (iv) o \textit{anti-reset windup} reduz o sobressinal, mas valores elevados de $K_{aw}$ esvaziam o integrador e deixam erro residual. O melhor compromisso foi obtido com o PIDF com $K_p = 3{,}5$, $K_i = 0{,}01$, $K_d = 1{,}5$ e $N = 0{,}05$. Ainda assim, nenhum ensaio de malha simples acomodou $L_2$ em menos de cerca de 55~s, e em todos eles o nível do tanque 1 continuou oscilando até o fim do ensaio, limitação que motiva o controle em cascata do Experimento 2C.
\end{justify}

\newpage

%=============================================================
\subsection{Experimento 2C: Controle em Cascata}

\begin{justify}
    \hspace{4ex}Nesta seção são apresentados os resultados do controle em cascata do nível do tanque 2. Em todas as figuras, a curva verde corresponde ao nível do tanque 1 ($L_1$), a curva vermelha ao nível do tanque 2 ($L_2$) e a linha horizontal em 15~cm à referência $r$. As métricas de todas as combinações, estimadas a partir dos gráficos conforme descrito na metodologia, estão reunidas na Tabela~\ref{tab:2c_combinacoes}, e os ganhos correspondentes, na Tabela~\ref{tab:2c_ganhos_metodologia}.
\end{justify}

\begin{justify}
    \hspace{4ex}Antes da análise individual, cabe destacar duas observações comuns a todos os ensaios. A primeira é que, no início de todos eles, $L_1$ cresce segundo uma rampa praticamente idêntica, o que indica que a bomba permanece saturada em $4$~V durante os primeiros 15 a 20~s, independentemente da combinação e dos ganhos. Assim, o tempo de subida de $L_2$ é limitado principalmente pela capacidade do atuador, e não pela sintonia, o que explica os valores semelhantes de $t_r$ (entre 15 e 20~s na maioria dos ensaios) na Tabela~\ref{tab:2c_combinacoes}; os valores menores, de cerca de 10~s, ocorreram nas combinações P/PI e P/PID, em que $L_2$ partiu de um nível mais alto e precisou percorrer uma variação menor. A segunda é que a bomba só consegue adicionar água ao tanque 2 indiretamente, via tanque 1, e nunca retirá-la: quando $L_2$ ultrapassa o valor final, o excesso só é eliminado pela vazão natural de saída do tanque 2, que é lenta. Por isso, qualquer sobressinal em $L_2$, mesmo pequeno, prolonga bastante o tempo de acomodação.
\end{justify}

%-------------------------------------------------------------
\subsubsection{Comparação entre Malha Simples e Controle em Cascata}

\begin{justify}
    \hspace{4ex}A Figura~\ref{fig:2c_simples_vs_cascata} compara o melhor resultado obtido com malha simples no Experimento 2B, o PIDF com $K_p = 3{,}5$, $K_i = 0{,}01$, $K_d = 1{,}5$ e $N = 0{,}05$, com a combinação PID/PID em cascata. A Tabela~\ref{tab:2c_simples_vs_cascata} resume as métricas, incluindo também o PID de malha simples com menor tempo de acomodação.
\end{justify}

\begin{figure}[H]
    \centering
    \subfloat[Malha simples -- PIDF]{\grafico{img/2b/PID-F/saida-p3.5-i0.01-d1.5-fn0.05.pdf}}\hfill
    \subfloat[Cascata -- mestre PID, escravo PID]{\grafico{img/2c/PID_PID/output.pdf}}
    \caption{Experimento 2C -- Comparação entre malha simples (Experimento 2B) e controle em cascata.}
    \label{fig:2c_simples_vs_cascata}
\end{figure}

\begin{table}[H]
\centering
\caption{Comparação entre malha simples e controle em cascata (valores aproximados).}
\label{tab:2c_simples_vs_cascata}
\resizebox{\textwidth}{!}{%
\begin{tabular}{|c|c|c|c|c|c|}
\hline
\textbf{Estratégia} & \textbf{Controlador(es) e ganhos} & $\mathbf{L_{1,max}}$ \textbf{(cm)} & $\mathbf{M_p}$ \textbf{(\%)} & $\mathbf{t_s}$ \textbf{(s)} & \textbf{Oscilação de} $\mathbf{L_1}$ \\ \hline
Malha simples & PID ($K_p = 3{,}5$, $K_i = 0{,}01$, $K_d = 0{,}8$) & 28 & 7 & 55 & até o fim do ensaio \\ \hline
Malha simples & PIDF ($K_p = 3{,}5$, $K_i = 0{,}01$, $K_d = 1{,}5$, $N = 0{,}05$) & 28 & 7 & 60 & até o fim do ensaio \\ \hline
Cascata & PID/PID (Tabela~\ref{tab:2c_ganhos_metodologia}) & 26 & 0 & 25 & não há \\ \hline
\end{tabular}%
}
\end{table}

\begin{justify}
    \hspace{4ex}A diferença entre as duas estratégias é evidente. Com uma única malha, o controlador enxerga a planta de segunda ordem completa: $L_1$ sobe até o patamar do intertravamento (em torno de 28~cm no gráfico), onde a tensão na bomba passa a ser limitada, e $L_2$ ultrapassa a referência em cerca de 7\%. Em seguida, $L_1$ passa a oscilar entre cerca de 12 e 18~cm e continua oscilando até o fim do ensaio, enquanto $L_2$ só se acomoda por volta de 60~s. Esse foi o melhor resultado de malha simples; com ganhos mais agressivos, como visto no Experimento 2B, o sobressinal chegou a 20\% e as oscilações não se extinguiram dentro dos 120~s.
\end{justify}

\begin{justify}
    \hspace{4ex}Com a cascata, por outro lado, $L_1$ sobe até cerca de 26~cm, sem tocar o intertravamento, e desce de forma suave e sem oscilações até o valor de regime; $L_2$ converge para a referência sem sobressinal visível e se acomoda em cerca de 25~s, menos da metade do tempo da malha simples e praticamente o mínimo imposto pela saturação da bomba. Isso ocorre porque a malha interna controla diretamente o nível do tanque 1, cuja dinâmica é mais rápida, e corrige os desvios de $L_1$ antes que eles se propaguem para o tanque 2. Como a malha interna é suficientemente mais rápida que a externa, o controlador mestre enxerga aproximadamente apenas a dinâmica do tanque 2, um sistema próximo de primeira ordem (Equação~\eqref{eq:ft_cascata}), o que facilita a sintonia e elimina as oscilações observadas com a malha simples. Além disso, como o mestre comanda uma referência de nível, e não uma tensão, a limitação de sua saída em 28~cm impõe naturalmente um teto para $L_1$, reduzindo a necessidade de atuação dos intertravamentos.
\end{justify}

%-------------------------------------------------------------
\subsubsection{Combinações de Controladores}

\begin{justify}
    \hspace{4ex}As Figuras~\ref{fig:2c_mestre_p} a~\ref{fig:2c_mestre_pidf} apresentam os resultados agrupados pelo controlador da malha mestre, e a Tabela~\ref{tab:2c_combinacoes} resume as métricas de todas as combinações.
\end{justify}

% ---------- Mestre P ----------
\paragraph{Mestre P.}

\begin{justify}
    \hspace{4ex}Com o mestre proporcional, todas as combinações apresentaram erro em regime de aproximadamente 2,5~cm em $L_2$, independentemente do controlador escravo (Figura~\ref{fig:2c_mestre_p}): o nível do tanque 2 estabiliza em torno de 12,5~cm, e não em 15~cm. Isso ocorre porque a malha externa não possui ação integral. Com um mestre puramente proporcional, a Equação~\eqref{eq:cascata_mestre} se reduz a $L_{1d} = K_{p,M}\,(r - L_2)$, em que $K_{p,M}$ é o ganho proporcional do mestre; assim, o mestre só pede ao tanque 1 um nível diferente de zero se houver erro em $L_2$. Como o tanque 1 precisa de um nível positivo para continuar alimentando o tanque 2, algum erro em $L_2$ sempre permanece.
\end{justify}

\begin{figure}[H]
    \centering
    \subfloat[Escravo P]{\grafico{img/2c/PP/output.pdf}}\hfill
    \subfloat[Escravo PI]{\grafico{img/2c/P_PI/output.pdf}}\\
    \subfloat[Escravo PID]{\grafico{img/2c/P_PID/output.pdf}}\hfill
    \subfloat[Escravo PIDF]{\grafico{img/2c/P_PIDF/OUTPUT.pdf}}
    \caption{Experimento 2C -- Controlador mestre P combinado com diferentes controladores escravos.}
    \label{fig:2c_mestre_p}
\end{figure}

\begin{justify}
    \hspace{4ex}O valor desse erro pode ser previsto a partir de três relações válidas em regime permanente. A primeira vem do próprio mestre P, conforme a Equação~\eqref{eq:2c_regime_mestre}.
\end{justify}

\begin{equation}
    L_{1d} = K_{p,M}\,(r - L_{2,ss})
    \label{eq:2c_regime_mestre}
\end{equation}

\begin{justify}
    \hspace{4ex}A segunda vem do escravo. Quando ele possui ação integral (PI, PID ou PIDF), seu integrador só para de variar quando o erro da malha interna é nulo, de modo que, em regime, o tanque 1 atinge exatamente o nível pedido pelo mestre, como indica a Equação~\eqref{eq:2c_regime_escravo}.
\end{justify}

\begin{equation}
    L_{1,ss} = L_{1d}
    \label{eq:2c_regime_escravo}
\end{equation}

\begin{justify}
    \hspace{4ex}A terceira vem da física da planta. Como discutido na Seção~\ref{sec:ft_nivel}, em regime as vazões de saída dos dois tanques se igualam e, pela lei de Torricelli, a razão entre os níveis depende apenas das áreas dos orifícios de saída. Escrevendo essa relação em termos do nível do tanque 1, obtém-se a Equação~\eqref{eq:2c_alpha}.
\end{justify}

\begin{equation}
    L_{1,ss} = \alpha\, L_{2,ss}, \qquad \alpha = \left(\frac{a_2}{a_1}\right)^2
    \label{eq:2c_alpha}
\end{equation}

\begin{justify}
    \hspace{4ex}O parâmetro $\alpha$ é, portanto, a razão entre os níveis dos dois tanques em regime, e depende apenas da relação entre as áreas dos orifícios de saída, e não dos controladores. Com orifícios iguais, $\alpha = 1$ e os dois tanques estabilizam no mesmo nível. Nos ensaios, $\alpha$ pode ser obtido diretamente dos gráficos, dividindo-se o nível final de $L_1$ pelo de $L_2$: os dois níveis terminam praticamente juntos ou com $L_1$ no máximo cerca de 1~cm acima de $L_2$, de modo que $\alpha$ fica entre aproximadamente 1 e 1,1.
\end{justify}

\begin{justify}
    \hspace{4ex}Substituindo as Equações~\eqref{eq:2c_regime_mestre} e~\eqref{eq:2c_regime_escravo} na Equação~\eqref{eq:2c_alpha}, obtém-se $K_{p,M}(r - L_{2,ss}) = \alpha\,L_{2,ss}$ e, isolando $L_{2,ss}$, chega-se à Equação~\eqref{eq:2c_ess_mestre_p}, que fornece o nível final do tanque 2 e, consequentemente, o erro em regime $e_{ss} = r - L_{2,ss}$.
\end{justify}

\begin{equation}
    L_{2,ss} = \frac{K_{p,M}}{K_{p,M} + \alpha}\, r
    \qquad\Rightarrow\qquad
    e_{ss} = \frac{\alpha}{K_{p,M} + \alpha}\, r
    \label{eq:2c_ess_mestre_p}
\end{equation}

\begin{justify}
    \hspace{4ex}A Equação~\eqref{eq:2c_ess_mestre_p} mostra que o erro só seria nulo com $K_{p,M}$ infinito e que ele não depende do controlador escravo, desde que este tenha ação integral. Para $r = 15$~cm, $K_{p,M} = 5$ e $\alpha \approx 1$, a equação prevê $L_{2,ss} \approx 12{,}5$~cm, isto é, $e_{ss} \approx 2{,}5$~cm, o que coincide com o observado nas combinações P/PID e P/PIDF. Para P/PI, com $K_{p,M} = 4$, o valor previsto é $L_{2,ss} = 12$~cm; no gráfico, $L_2$ ainda está descendo lentamente ao final do ensaio, aproximando-se desse valor. A equação também indica que, para reduzir o erro a menos de 1~cm, seria necessário $K_{p,M} \approx 14$, o que tornaria a malha externa muito agressiva. A solução adequada é incluir a ação integral no mestre. No caso P/P, como o escravo não tem integrador, a Equação~\eqref{eq:2c_regime_escravo} não vale e o tanque 1 fica um pouco abaixo de $L_{1d}$, mas o erro em $L_2$ observado é da mesma ordem, cerca de 2,5~cm.
\end{justify}

\begin{justify}
    \hspace{4ex}O controlador escravo, embora não altere o erro em regime, altera bastante o transitório. Com o escravo P (P/P), a resposta de $L_2$ é sobreamortecida e acomoda em cerca de 20~s, pois o escravo deixa de saturar mais cedo e $L_1$ atinge o menor pico entre todos os ensaios (cerca de 23~cm). Já com os escravos com ação integral, sintonizados com $K_i$ elevado ($0{,}2$ e $0{,}175$), o integrador do escravo acumula erro durante todo o período de saturação da bomba, e $L_1$ sobe até cerca de 26,5~cm. Esse excesso de água no tanque 1 é transferido ao tanque 2, que sobe até cerca de 14,5~cm, aproximadamente 15\% acima do valor final, e depois desce lentamente, apenas pela vazão de saída, levando cerca de 65 a 70~s para acomodar. Com o escravo PIDF, o efeito é ainda mais intenso: $L_1$ atinge o limite de 28~cm imposto à saída do mestre, onde permanece por alguns segundos, e $L_2$ chega a cerca de 16~cm, aproximadamente 30\% acima do valor final. Portanto, com mestre P, a ação integral no escravo não traz benefício em regime e piora o transitório.
\end{justify}

% ---------- Mestre PI ----------
\paragraph{Mestre PI.}

\begin{justify}
    \hspace{4ex}Com ação integral na malha mestre, o erro em regime de $L_2$ foi eliminado em todas as combinações (Figura~\ref{fig:2c_mestre_pi}). Isso pode ser entendido pela mesma análise de regime feita para o mestre P: o integrador do mestre só para de variar quando $e_x = r - L_2 = 0$, de modo que ele ajusta $L_{1d}$ até o valor $\alpha\, r$ necessário para que $L_2 = r$, sem que seja preciso manter erro algum. As combinações PI/P, PI/PI e PI/PID tiveram respostas praticamente idênticas: $L_2$ converge para 15~cm sem sobressinal visível, com tempo de acomodação entre 25 e 30~s, e $L_1$ atinge um pico de cerca de 26~cm antes de descer suavemente até o valor de regime.
\end{justify}

\begin{figure}[H]
    \centering
    \subfloat[Escravo P]{\grafico{img/2c/PI_P/output.pdf}}\hfill
    \subfloat[Escravo PI]{\grafico{img/2c/PI_PI/output.pdf}}\\
    \subfloat[Escravo PID]{\grafico{img/2c/PI_PID/output.pdf}}\hfill
    \subfloat[Escravo PIDF]{\grafico{img/2c/PI_PIDF/OUTPUT.pdf}}
    \caption{Experimento 2C -- Controlador mestre PI combinado com diferentes controladores escravos.}
    \label{fig:2c_mestre_pi}
\end{figure}

\begin{justify}
    \hspace{4ex}O resultado da combinação PI/P mostra que, com o integrador no mestre, o escravo apenas proporcional já é suficiente. Nesse caso, a Equação~\eqref{eq:2c_regime_escravo} deixa de valer e o tanque 1 fica um pouco abaixo de $L_{1d}$, mas esse erro da malha interna é compensado pelo integrador do mestre, que eleva $L_{1d}$ até que $L_2 = r$. A inclusão da ação integral no escravo (PI/PI), com ganho baixo ($K_i = 0{,}03$), não trouxe ganho perceptível de desempenho, mas também não causou oscilações. Da mesma forma, a ação derivativa do escravo PID ($K_d = 0{,}01$) não alterou a resposta, como será discutido na Seção~\ref{sec:2c_ganhos}.
\end{justify}

\begin{justify}
    \hspace{4ex}Na combinação PI/PIDF, $L_1$ chegou ao limite de 28~cm da saída do mestre e $L_2$ ultrapassou ligeiramente a referência, chegando a cerca de 15,5~cm (sobressinal da ordem de 3\%), e levou aproximadamente 65~s para acomodar. Como o filtro atua apenas sobre uma ação derivativa de ganho muito pequeno, esse sobressinal não se deve ao filtro, e sim à saturação da saída do mestre: enquanto $L_{1d}$ está limitado em 28~cm, o integrador do mestre, que não possui \textit{anti-windup}, continua acumulando erro; quando $L_2$ se aproxima da referência, esse excesso mantém $L_{1d}$ elevado por mais tempo e, como o tanque 2 não pode ser esvaziado pela bomba, o sobressinal resultante decai lentamente. Observa-se ainda uma pequena perturbação em $L_1$ por volta de 60~s, que é corrigida pela malha interna sem efeito visível em $L_2$, exemplificando a capacidade da cascata de rejeitar distúrbios no tanque 1.
\end{justify}

% ---------- Mestre PID ----------
\paragraph{Mestre PID.}

\begin{justify}
    \hspace{4ex}Os resultados com mestre PID (Figura~\ref{fig:2c_mestre_pid}) são muito próximos dos obtidos com mestre PI. As combinações PID/P, PID/PI e PID/PID levaram $L_2$ à referência sem sobressinal visível e com tempo de acomodação entre 25 e 30~s; as diferenças em relação ao mestre PI são de poucos segundos, comparáveis à repetibilidade dos ensaios. A ação derivativa do mestre, com $K_d = 0{,}05$, tem efeito muito pequeno, pois $L_2$ varia lentamente.
\end{justify}

\begin{figure}[H]
    \centering
    \subfloat[Escravo P]{\grafico{img/2c/PID_P/output.pdf}}\hfill
    \subfloat[Escravo PI]{\grafico{img/2c/PID_PI/output.pdf}}\\
    \subfloat[Escravo PID]{\grafico{img/2c/PID_PID/output.pdf}}\hfill
    \subfloat[Escravo PIDF]{\grafico{img/2c/PID_PIDF/OUTPUT.pdf}}\\
    \subfloat[Escravo PIDAW]{\grafico{img/2c/PID_PIDAW/output.pdf}}
    \caption{Experimento 2C -- Controlador mestre PID combinado com diferentes controladores escravos.}
    \label{fig:2c_mestre_pid}
\end{figure}

\begin{justify}
    \hspace{4ex}As combinações PID/PIDF e PID/PIDAW, realizadas com os mesmos ganhos proporcional, integral e derivativo no escravo, permitem isolar o efeito do \textit{anti-reset windup}. Com o PIDF no escravo, o integrador continua acumulando erro enquanto a bomba está saturada; $L_1$ chega ao limite de 28~cm, $L_2$ ultrapassa a referência em cerca de 0,5~cm (aproximadamente 3\%) e só acomoda por volta de 70~s. Com o PIDAW, o integrador é descarregado durante a saturação (Equação~\eqref{eq:ft_aw}), o que reduz o pico de $L_1$ para cerca de 27,5~cm, praticamente elimina o sobressinal de $L_2$ e reduz o tempo de acomodação para cerca de 35~s, aproximadamente metade do valor anterior. Esse resultado confirma que, nesta planta, o \textit{windup} do escravo durante a saturação inicial é a principal causa do sobressinal em $L_2$.
\end{justify}

% ---------- Mestre PIDF ----------
\paragraph{Mestre PID com filtro derivativo (PIDF).}

\begin{justify}
    \hspace{4ex}Além das combinações exigidas no roteiro, o mestre PID com filtro derivativo também foi testado (Figura~\ref{fig:2c_mestre_pidf}). Com os escravos P, PI, PID e PIDF, o comportamento foi muito semelhante: $L_1$ chega ao limite de 28~cm da saída do mestre, $L_2$ ultrapassa a referência em cerca de 0,5~cm (sobressinal da ordem de 3\%) e o tempo de acomodação fica entre 50 e 65~s, sem erro em regime. Assim como nos casos anteriores, a escolha do escravo entre P, PI, PID e PIDF teve pouca influência, pois a dinâmica é dominada pela saturação inicial da bomba e pela ação integral do mestre. O filtro no mestre também não produziu diferença visível no nível: com $N = 0{,}005$~rad/s, a frequência de corte da Equação~\eqref{eq:ft_filtro} é muito inferior às frequências da resposta, e a parcela derivativa fica praticamente anulada; seu benefício esperado, a redução do ruído transmitido a $L_{1d}$ e a $V_P$, não pode ser verificado apenas pelas curvas de nível.
\end{justify}

\begin{figure}[H]
    \centering
    \subfloat[Escravo P]{\grafico{img/2c/PIDF_P/output.pdf}}\hfill
    \subfloat[Escravo PI]{\grafico{img/2c/PIDF_PI/output.pdf}}\\
    \subfloat[Escravo PID]{\grafico{img/2c/PIDF_PID/output.pdf}}\hfill
    \subfloat[Escravo PIDF]{\grafico{img/2c/PIDF_PIDF/OUTPUT.pdf}}\\
    \subfloat[Escravo PIDAW]{\grafico{img/2c/PIDF_PIDAW/OUTPUT.pdf}}
    \caption{Experimento 2C -- Controlador mestre PID com filtro derivativo combinado com diferentes controladores escravos.}
    \label{fig:2c_mestre_pidf}
\end{figure}

\begin{justify}
    \hspace{4ex}A combinação PIDF/PIDAW foi a de melhor desempenho entre as combinações com mestre PIDF: $L_1$ teve o menor pico do grupo (cerca de 27~cm), e $L_2$ alcançou a referência sem sobressinal visível e acomodou em cerca de 30~s. Comparada à combinação PIDF/PIDF, o tempo de acomodação caiu para menos da metade (de cerca de 65 para 30~s), reforçando a conclusão obtida com o mestre PID: o \textit{anti-reset windup} na malha escrava é o elemento que mais melhora a resposta, porque atua justamente no período em que a bomba está saturada. Vale notar que, nesse ensaio, o escravo usou $K_i = 0{,}05$, maior que nos demais, e mesmo assim não houve sobressinal, já que o \textit{anti-windup} impede que esse ganho maior se traduza em acúmulo excessivo durante a saturação.
\end{justify}

\begin{table}[H]
\centering
\caption{Métricas aproximadas das combinações de controladores no controle em cascata, estimadas a partir dos gráficos. $L_{1,\max}$: pico de $L_1$; $L_{1,ss}$: nível de regime de $L_1$. $M_p$ calculado em relação ao valor final de $L_2$, conforme a Equação~\eqref{eq:ft_mp}.}
\label{tab:2c_combinacoes}
\resizebox{\textwidth}{!}{%
\begin{tabular}{|c|c|c|c|c|c|c|c|c|}
\hline
\textbf{Mestre} & \textbf{Escravo} & $\mathbf{L_{1,max}}$ \textbf{(cm)} & $\mathbf{L_{2,max}}$ \textbf{(cm)} & $\mathbf{M_p}$ \textbf{(\%)} & $\mathbf{t_r}$ \textbf{(s)} & $\mathbf{t_s}$ \textbf{(s)} & $\mathbf{e_{ss}}$ \textbf{(cm)} & $\mathbf{L_{1,ss}}$ \textbf{(cm)} \\ \hline
P    & P     & 23   & 12,5 & 0   & 15 & 20 & 2,5 & 12   \\ \hline
P    & PI    & 26,5 & 14,5 & 15  & 10 & 65 & 2,5 & 12   \\ \hline
P    & PID   & 26,5 & 14,5 & 15  & 10 & 70 & 2,5 & 12   \\ \hline
P    & PIDF  & 28   & 16   & 30  & 15 & 75 & 2,5 & 13,5 \\ \hline
PI   & P     & 26   & 15   & 0   & 15 & 30 & 0   & 15   \\ \hline
PI   & PI    & 26   & 15   & 0   & 15 & 25 & 0   & 15   \\ \hline
PI   & PID   & 26   & 15   & 0   & 15 & 25 & 0   & 15   \\ \hline
PI   & PIDF  & 28   & 15,5 & 3   & 20 & 65 & 0   & 16   \\ \hline
PID  & P     & 26   & 15   & 0   & 15 & 30 & 0   & 15   \\ \hline
PID  & PI    & 26   & 15   & 0   & 15 & 25 & 0   & 15   \\ \hline
PID  & PID   & 26   & 15   & 0   & 15 & 25 & 0   & 15   \\ \hline
PID  & PIDF  & 28   & 15,5 & 3   & 20 & 70 & 0   & 16   \\ \hline
PID  & PIDAW & 27,5 & 15   & 0   & 20 & 35 & 0   & 16   \\ \hline
PIDF & P     & 28   & 15,5 & 3   & 20 & 50 & 0   & 16,5 \\ \hline
PIDF & PI    & 28   & 15,5 & 3   & 20 & 55 & 0   & 16   \\ \hline
PIDF & PID   & 28   & 15,5 & 3   & 20 & 55 & 0   & 16   \\ \hline
PIDF & PIDF  & 28   & 15,5 & 3   & 15 & 65 & 0   & 16   \\ \hline
PIDF & PIDAW & 27   & 15   & 0   & 20 & 30 & 0   & 16   \\ \hline
\end{tabular}%
}
\end{table}

%-------------------------------------------------------------
\subsubsection{Influência dos Ganhos}
\label{sec:2c_ganhos}

\begin{justify}
    \hspace{4ex}Para cada combinação foi registrado apenas o ensaio com o conjunto final de ganhos (Tabela~\ref{tab:2c_ganhos_metodologia}); entretanto, como os ganhos variam entre as combinações, é possível avaliar a influência de cada um deles comparando ensaios semelhantes.
\end{justify}

\begin{itemize}
    \item \textbf{Ganho proporcional do mestre ($K_{p,M}$):} com mestre P, determina diretamente o erro em regime, conforme a Equação~\eqref{eq:2c_ess_mestre_p}. Passar de $K_{p,M} = 4$ para $5$ reduz o erro previsto de cerca de 3 para 2,5~cm, uma melhoria pequena, o que evidencia que o mestre P não é adequado quando se exige erro nulo. Com mestre PI ou PID, $K_{p,M}$ entre 4 e 4,5 produziu respostas sem oscilação, o que indica que a malha interna é rápida o suficiente para que o mestre opere com ganhos relativamente altos.

    \item \textbf{Ganho integral do mestre ($K_{i,M}$):} é o responsável pela eliminação do erro em regime. Valores entre 0,07 e 0,1 eliminaram o erro em todos os casos. Nos ensaios em que $L_{1d}$ atingiu o limite de 28~cm, porém, o integrador do mestre continuou acumulando erro com a saída saturada, o que produziu um pequeno sobressinal em $L_2$, de cerca de 0,5~cm. Um \textit{anti-windup} também na malha mestre poderia reduzir esse efeito.

    \item \textbf{Ganho integral do escravo ($K_{i,E}$):} quando elevado ($0{,}175$ a $0{,}2$, nas combinações P/PI, P/PID e P/PIDF), causou acúmulo de erro durante a saturação da bomba, elevando o pico de $L_1$ e o sobressinal de $L_2$. Com valores baixos ($0{,}03$), o escravo PI ou PID teve desempenho equivalente ao escravo P. Na prática, quando o mestre já possui ação integral, o integrador do escravo é dispensável; quando é usado, deve vir acompanhado de \textit{anti-windup}.

    \item \textbf{Ganhos derivativos ($K_d$) e filtro ($N$):} com os valores utilizados, a contribuição da ação derivativa é desprezível. Pelos gráficos, durante a subida $L_2$ varia cerca de 10~cm em 15~s, ou seja, aproximadamente $0{,}7$~cm/s. Pela Equação~\eqref{eq:ft_pid}, a parcela derivativa do mestre é $K_{d,M}\, de_x/dt$, e, como a referência é constante, $de_x/dt = -dL_2/dt$; com $K_{d,M} = 0{,}05$, essa parcela vale apenas cerca de $0{,}05 \times 0{,}7 \approx 0{,}04$~cm em $L_{1d}$. Da mesma forma, com $L_1$ variando cerca de 1~cm/s, a parcela derivativa do escravo, com $K_{d,E} = 0{,}01$, é da ordem de $0{,}01$~V em $V_P$, desprezível diante dos 4~V de saturação. Por isso, as combinações com e sem ação derivativa (PI/PI e PI/PID, PID/PI e PID/PID) apresentaram curvas praticamente sobrepostas. Nos controladores com filtro, o valor $N = 0{,}005$~rad/s reduz ainda mais essa parcela, pois atenua praticamente todas as frequências presentes na resposta. Ganhos derivativos maiores não foram necessários, pois a cascata já eliminou as oscilações que, na malha simples, justificavam o uso da ação derivativa.

    \item \textbf{Ganho de \textit{anti-windup} ($K_{aw}$):} foi o parâmetro de maior impacto entre os ensaios comparáveis. Pela Equação~\eqref{eq:ft_aw}, enquanto a bomba está saturada, o termo $K_{aw}(u_{sat} - u)$ é negativo e descarrega o integrador. Com isso, o tempo de acomodação caiu para aproximadamente metade (PID/PIDF $\rightarrow$ PID/PIDAW) e para menos da metade (PIDF/PIDF $\rightarrow$ PIDF/PIDAW), e o sobressinal foi eliminado.
\end{itemize}

\begin{justify}
    \hspace{4ex}De forma geral, os resultados confirmam a regra de sintonia de sistemas em cascata: a malha interna deve ser mais rápida que a externa e deve ser projetada primeiro. Nesta planta, como o tanque 1 já é naturalmente mais rápido que o tanque 2 e a bomba permanece saturada no início, o desempenho foi pouco sensível aos ganhos proporcionais e derivativos e muito sensível ao tratamento da saturação, isto é, à quantidade de erro acumulada pelos integradores enquanto a bomba está saturada.
\end{justify}

%-------------------------------------------------------------
\subsubsection{Síntese do Experimento 2C}

\begin{justify}
    \hspace{4ex}O controle em cascata superou com folga o controle com malha simples do Experimento 2B: eliminou as oscilações persistentes de $L_1$ e o sobressinal de $L_2$, manteve $L_1$ limitado a no máximo 28~cm pela saturação da saída do mestre, sem oscilar contra o patamar do intertravamento, e reduziu o tempo de acomodação de cerca de 55 a 60~s, nos melhores ensaios de malha simples, para cerca de 25~s. Entre as combinações, a presença de ação integral no mestre é o requisito essencial para erro nulo; o controlador escravo pode ser simples (P), desde que não acumule erro durante a saturação. As melhores combinações foram PIDF/PIDAW (cerca de 30~s, sem sobressinal) e PI/PI, PI/PID, PID/PI e PID/PID (cerca de 25~s, sem sobressinal). Considerando também a simplicidade de implementação, a combinação PI/P é uma alternativa adequada, com desempenho muito próximo (cerca de 30~s, sem sobressinal) e apenas três ganhos a ajustar.
\end{justify}

\newpage

%%%%%%%%%% CONCLUSÃO %%%%%%%%%%%%%%%

\thispagestyle{main}

\section{CONCLUSÃO}

\begin{justify}
    \hspace{4ex}Os experimentos realizados permitiram analisar o comportamento de diferentes estratégias de controle aplicadas ao sistema de tanques acoplados da Quanser, relacionando os conceitos teóricos de controle em malha fechada aos resultados obtidos experimentalmente. A implementação dos controladores no ambiente MATLAB/SIMULINK possibilitou avaliar a influência das ações proporcional, integral e derivativa, bem como os efeitos da saturação do atuador, da sintonia dos ganhos e da utilização de mecanismos de compensação sobre o desempenho dinâmico da planta.
\end{justify}

\begin{justify}
    \hspace{4ex}No Experimento 2A, referente ao controle do nível do tanque 1, observou-se que a ação proporcional isolada não foi suficiente para eliminar o erro em regime permanente, enquanto a inclusão da ação integral permitiu atingir a referência desejada. Entretanto, ganhos integrais elevados provocaram sobressinal e aumento do tempo de acomodação, evidenciando a importância de uma sintonia adequada. Os melhores resultados foram obtidos com os controladores PID, PIDF e PIDAW, que apresentaram respostas com tempo de acomodação de aproximadamente 10~s, sem sobressinal significativo e com erro em regime praticamente nulo. Dessa forma, verificou-se que a utilização de estruturas mais complexas não necessariamente proporciona melhorias expressivas quando um controlador mais simples já se encontra adequadamente ajustado.
\end{justify}

\begin{justify}
    \hspace{4ex}No Experimento 2B, o controle do nível do tanque 2 apresentou maiores dificuldades devido à dinâmica de segunda ordem do sistema, caracterizada pelo acoplamento entre os tanques. O controlador P manteve erro em regime e já apresentou resposta subamortecida com ganho baixo, e ganhos integrais elevados levaram a \textit{windup} severo, com sobressinal de até cerca de 50\% e oscilações sem acomodação. A ação derivativa reduziu o sobressinal e o tempo de acomodação, e o filtro eliminou o ruído que ela introduzia no sinal de controle; o melhor compromisso foi obtido com o PIDF, com sobressinal de cerca de 7\% e acomodação em aproximadamente 60~s. O \textit{anti-reset windup} reduziu o sobressinal, mas valores elevados de $K_{aw}$ esvaziaram o integrador e deixaram erro residual. Em todos os ensaios, porém, o nível do tanque 1 continuou oscilando até o fim do ensaio. Esses resultados evidenciaram as limitações do controle em malha simples diante da dinâmica mais lenta da planta e das restrições físicas impostas à atuação da bomba.
\end{justify}

\begin{justify}
    \hspace{4ex}No Experimento 2C, a implementação do controle em cascata proporcionou uma melhoria significativa no desempenho do sistema. A utilização de uma malha interna para controlar o nível do tanque 1 permitiu corrigir mais rapidamente suas variações, reduzindo os efeitos propagados ao tanque 2. Na comparação entre o melhor resultado de malha simples e a combinação PID/PID em cascata, o tempo de acomodação foi reduzido de cerca de 60~s para aproximadamente 25~s, com eliminação das oscilações persistentes do tanque 1 e ausência de sobressinal visível. Verificou-se também que a presença da ação integral no controlador mestre foi fundamental para eliminar o erro em regime permanente, enquanto a ação derivativa apresentou influência reduzida nas condições e nos ganhos utilizados. Entre as combinações avaliadas, destacaram-se PI/PI, PI/PID, PID/PI e PID/PID, com tempos de acomodação próximos de 25~s, além da combinação PIDF/PIDAW, que apresentou acomodação em aproximadamente 30~s, sem sobressinal visível, demonstrando a importância do mecanismo de \textit{anti-reset windup} na redução dos efeitos da saturação.
\end{justify}

\begin{justify}
    \hspace{4ex}Por fim, os resultados demonstraram que a eficiência de um sistema de controle depende não apenas da escolha do controlador, mas também da sintonia de seus parâmetros, da arquitetura adotada e do respeito às limitações físicas da planta. A utilização de intertravamentos e da saturação do sinal de controle constitui um aspecto importante para a segurança operacional, enquanto o tratamento adequado da ação integral contribui para evitar respostas transitórias excessivas. Assim, o laboratório permitiu consolidar os conhecimentos sobre controladores PID e suas variações, evidenciando as vantagens do controle em cascata para sistemas com dinâmicas acopladas e reforçando a importância de conciliar desempenho, simplicidade de implementação e segurança na aplicação prática de sistemas de controle.
\end{justify}

\newpage

%%%%%%%% REFERÊNCIAS %%%%%%%%%%%%%%%%%

%\newpage
\phantomsection
\addcontentsline{toc}{section}{REFERÊNCIAS}

\bibliographystyle{ieeetr}
\bibliography{bib/bibliografia}

\end{document}

